\documentclass[11pt, a4paper]{report} % add draft if needed

\usepackage[english]{babel}
\usepackage[T1]{fontenc}
\usepackage[utf8]{inputenc}
\usepackage{mathtools}
\usepackage{amsfonts}
\usepackage{listings}
\usepackage[a4paper]{geometry}
\usepackage[parfill]{parskip}
\usepackage{hyperref}
\usepackage{lipsum}
\usepackage[toc,page]{appendix}
\usepackage{lastpage}
\usepackage{fancyhdr}
\usepackage[explicit]{titlesec}
\usepackage[style=alphabetic]{biblatex}
\usepackage{color}
\usepackage{amsthm}
\usepackage{todo}
\usepackage{tikz}
\usepackage{subcaption}
\usepackage{caption}
\usepackage{chemfig}
\usepackage{cleveref}
\usepackage{multicol}
\usepackage{minted}
\usepackage{mdframed}

\providecommand{\tightlist}{%
  \setlength{\itemsep}{0pt}\setlength{\parskip}{0pt}}

\setminted{linenos,frame=lines,framesep=2mm}
\setminted[text]{linenos=false}


%tikz extras
\usetikzlibrary{arrows, shapes, snakes, automata, backgrounds, positioning}

\tikzset{
    triple/.style args={[#1] in [#2] in [#3]}{
        #1,preaction={preaction={draw,#3},draw,#2}
    }
}

% math envirs
\theoremstyle{definition}
\newtheorem{defn}{Definition}
\newtheorem{obsv}{Observation}

\newenvironment{observation}
    {\begin{framed}\begin{obsv}}
    {\end{obsv}\end{framed}}

% Chapter title formatting
\definecolor{gray75}{gray}{0.75}
\newcommand{\hsp}{\hspace{20pt}}

\titleformat{\chapter}[display]{}{
}{-85pt}{\Huge\bfseries\centering\textcolor{gray75}{|}\hsp#1\hsp\textcolor{gray75}{|}\\\hspace{10pt}}

% Header and footer
\fancypagestyle{plain}{% plain style is used on chapter pages
  \fancyhf{}%
  \fancyfoot[C]{\thepage}% Page \thepage\ of \pageref*{LastPage}
  \renewcommand{\headrulewidth}{0pt}% Line at the header invisible
  %\renewcommand{\footrulewidth}{0.4pt}% Line at the footer visible
}

\pagestyle{fancy}
\fancypagestyle{MyStyle}{
    \fancyhf{}
    \fancyhead[L]{ \leftmark }
    \fancyhead[R]{ \today }
    \fancyfoot[C]{ \thepage } % chktex 2
    \renewcommand{\chaptermark}[1]{\markboth{\textsl{##1}}{}}
    \renewcommand{\headrulewidth}{0pt}
}

\newcommand{\blankpage}{
  \newpage
  \thispagestyle{empty}
  \mbox{}
  \newpage
}

% \definecolor{mygreen}{rgb}{0,0.6,0}
% \definecolor{mygray}{rgb}{0.5,0.5,0.5}
% \definecolor{mymauve}{rgb}{0.58,0,0.82}

% \lstset{ %
%   backgroundcolor=\color{white},   % choose the background color; you must add \usepackage{color} or \usepackage{xcolor}; should come as last argument
%   basicstyle=\footnotesize,        % the size of the fonts that are used for the code
%   breakatwhitespace=false,         % sets if automatic breaks should only happen at whitespace
%   breaklines=true,                 % sets automatic line breaking
%   captionpos=b,                    % sets the caption-position to bottom
%   commentstyle=\color{mygreen},    % comment style
%   deletekeywords={...},            % if you want to delete keywords from the given language
%   escapeinside={\%*}{*)},          % if you want to add LaTeX within your code
%   extendedchars=true,              % lets you use non-ASCII characters; for 8-bits encodings only, does not work with UTF-8
%   frame=single,                       % adds a frame around the code
%   keepspaces=true,                 % keeps spaces in text, useful for keeping indentation of code (possibly needs columns=flexible)
%   keywordstyle=\color{blue},       % keyword style
%   language=Octave,                 % the language of the code
%   morekeywords={*,...},           % if you want to add more keywords to the set
%   numbers=left,                    % where to put the line-numbers; possible values are (none, left, right)
%   numbersep=5pt,                   % how far the line-numbers are from the code
%   numberstyle=\tiny\color{mygray}, % the style that is used for the line-numbers
%   rulecolor=\color{black},         % if not set, the frame-color may be changed on line-breaks within not-black text (e.g. comments (green here))
%   showspaces=false,                % show spaces everywhere adding particular underscores; it overrides 'showstringspaces'
%   showstringspaces=false,          % underline spaces within strings only
%   showtabs=false,                  % show tabs within strings adding particular underscores
%   stepnumber=2,                    % the step between two line-numbers. If it's 1, each line will be numbered
%   stringstyle=\color{mymauve},     % string literal style
%   tabsize=2,                       % sets default tabsize to 2 spaces
%   title=\lstname                   % show the filename of files included with \lstinputlisting; also try caption instead of title
% }

\newcommand{\joliel}[1]{\mintinline{jolie}{#1}}
\newcommand{\txtl}[1]{\mintinline{text}{#1}}

% Macros for component references
\newcommand{\cache}{\txtl{cache} }
\newcommand{\registry}{\txtl{registry} }
\newcommand{\regdb}{\txtl{registry-db} }
\newcommand{\security}{\txtl{security} }
\newcommand{\bcrypt}{\txtl{bcrypt} }

\title{Building a Package Manager for Jolie}
\author{Dan Sebastian Thrane}

\addbibresource{biblio.bib}

\begin{document}
\begin{titlepage}

\newcommand{\HRule}{\rule{\linewidth}{0.5mm}} % Defines a new command for the horizontal lines, change thickness here

%\ % Center everything on the page
\flushright %centering
%----------------------------------------------------------------------------------------
%   LOGO SECTION
%----------------------------------------------------------------------------------------
%\begin{minipage}{0.9\textwidth}
%\centering
%\hspace*{2cm}
%
\includegraphics[scale=0.35,bb=0 0 250 0]{pictures/SDU_logo.eps}\\[1cm] % Include a department/university logo - this will require the graphicx package

%----------------------------------------------------------------------------------------
%   HEADING SECTIONS
%----------------------------------------------------------------------------------------

\textsc{\LARGE University of Southern Denmark}\\%[0.3cm] % Name of your university/college
\textsc{\normalsize Department of Mathematics and Computer Science,\\ IMADA}\\[2.6cm]
 % Major heading such as course name
 % Minor heading such as course title

%----------------------------------------------------------------------------------------
%   TITLE SECTION
%----------------------------------------------------------------------------------------
\textsc{\Large Master Thesis}\\[0.5cm]
\hfill \\%[0.10cm]
{ \Huge \bfseries Building a Package Manager for Jolie }\\[0.4cm] % Title of your document
\hfill \\[2.8cm]

%----------------------------------------------------------------------------------------
%   AUTHOR SECTION
%----------------------------------------------------------------------------------------
\large
\emph{Author:}\\
Dan Sebastian \textsc{Thrane} % Your name
\\[0.8cm]
%~
\large
\emph{Supervisor:} \\
Fabrizio \textsc{Montesi} % Supervisor's Name


\vfill
% If you don't want a supervisor, uncomment the two lines below and remove the section above
%\Large \emph{Author:}\\
%John \textsc{Smith}\\[3cm] % Your name

%----------------------------------------------------------------------------------------
%   DATE SECTION
%----------------------------------------------------------------------------------------

{\large \today}%\\[2cm] % Date, change the \today to a set date if you want to be precise
%\end{minipage}

\end{titlepage}



\tableofcontents

\blankpage

\pagestyle{MyStyle}

\chapter{Introduction}
TODO Something which gets us onto the package track.

Packages come in many different forms, and what they are mostly depends on the
system that they serve. For example, binary package managers may ship
applications which are then installed on the user’s system. An example of
such as package manager could be the iOS App Store. In general, a package
manager primarily makes several things easier:

\begin{enumerate}
    \item Publishing of Packages
    \item Find packages to install
    \item Automatically install dependencies
    \item Handle available updates
\end{enumerate}

These same principles have also been applied to the application level.
Application-level packages are usually software components, which are intended
to be run locally as a part of your application. Package managers also
typically attempt to establish some sort of language/framework convention for
developing packages. This is done such that packages may be imported into a
software project in a consistent manner.

The goal of this thesis is to develop a package manager for the
Jolie\autocite{montesi2010jolie,JOLA} ecosystem. Jolie is a service-oriented
programming language.
% Do we really need this?
As such, we aim not only at covering the local packages (i.e., services running
        locally), but also external services. From the perspective of a local
Jolie service, an external service is any service, which runs on an external
virtual machine.

Before being able to achieve the goal of creating a package manager, a few
hurdles must first be overcome. In this thesis, we will present a new natively
supported module and configuration system for Jolie.

% Should also cover the somewhat blurry line there is between package manager
% and build tool

% Method
% ------
% This section focuses a bit too much on looking at other products. We should
% still mention a few up here. But we never actually looked straight at other
% products.
%
% There is also no focus on the configuration system which is developed.

\pagebreak

\section{Structure of this Thesis}

\textbf{Chapter 1}. This chapter. Introduces the thesis, explains the structure
of it, and notations used throughout the thesis.

\textbf{Chapter 2}. Covers background material which is relevant for the entire
thesis.

\textbf{Chapter 3}. Covers the module system and the configuration system which
had been made to support the package manager efforts.

\textbf{Chapter 4}. Explains the concept of ``Jolie Packages'' which are what
the Jolie Package Manager (JPM) supports. This chapter will also cover many
of the features that are directly related to packages, which JPM implements.

\textbf{Chapter 5}. Covers architectural and implementation specific details
for JPM.

\textbf{Chapter 6}. Conclusion of this thesis.

\section{Notation Used}

\subsection*{Service Diagrams}

Simple box-and-line diagrams are used for illustrating how service are
deployed, and how they depend on eachother. Figure
\ref{fig:simple_box_and_line} shows a simple diagram with two services \txtl{A}
and \txtl{B}. In this diagram \txtl{A} depends on the services provided by
\txtl{B}.

\begin{figure}[H]

\centering
\includegraphics[width=0.5\textwidth]{introduction/simple_box_and_line.eps}

\caption{A service diagram showing two services: \txtl{A} and \txtl{B}. In this
    diagram service \txtl{A} depends on service \txtl{B}}

\label{fig:simple_box_and_line}

\end{figure}

We'll illustrate embedding (a Jolie concept, see Chapter \ref{cha:background})
by drawing a service inside another. The parent service will show its
dependency on this service through similar means.

\begin{figure}[H]

\centering
\includegraphics[width=0.5\textwidth]{introduction/embedding.eps}

\caption{A \txtl{parent} service embedding a \txtl{child} service}

\label{fig:box_and_line_embedding}

\end{figure}

\subsection*{File Listings}

When code snippets are presented, where the file name and location is relevant
then the file name and file path will be shown in a comment before the file
contents. We will commonly show multiple files in a single listing.  As an
example, if we have the directory structure shown in Listing
\ref{lst:file_structure_example}.

\begin{listing}[H]
\begin{minted}{text}
/
|-- a
|   `-- 1
|-- b
|   `-- 2
`-- c
    `-- 3

3 directories, 3 files
\end{minted}
\caption{We will illustrate a directory structure in this manner}
\label{lst:file_structure_example}
\end{listing}

And the file contents of file \txtl{/a/1} is ``A'', \txtl{/b/2} contains ``B'',
    and \txtl{/c/3} contains ``C'', then we would show the contents of these
    files as in Listing \ref{lst:multiple_files_example}.

\begin{listing}[H]
\begin{minted}{text}
/* /a/1 */
A

/* /b/2 */
B

/* /c/3 */
C
\end{minted}

\caption{Multiple files shown in a single listing where file path and name is
    relevant.}

\label{lst:multiple_files_example}
\end{listing}

\subsection*{Command Line Illustrations}

In certain cases, examples are made by showing the output of a command line
session. The reader is expected to be familiar with a basic Unix-like shell.
When the working directory is relevant, then it will be shown as part of the
prompt. An example of this is shown in Listing \ref{lst:shell_example}.  In
this example two commands are invoked (\txtl{ls} and \txtl{cat hello-world.ol}.
The \verb!#! symbols mark the beginning of a command. Lines not
containing this symbol should be considered as part of the output from
the previous command. Before that the absolute path to the current
working directory is shown, in these cases they were both
\txtl{/this/is/the/path}.

\begin{listing}[H]
\begin{minted}{text}
dan@host:/this/is/the/path # ls
hello-world.ol

dan@host:/this/is/the/path # cat hello-world.ol
include "console.iol"
main {
    println@Console("Hello, World!")()
}
\end{minted}
\caption{Illustration of two commands being invoked in a shell}
\label{lst:shell_example}
\end{listing}

\section{Implementation Status}



\chapter{Background}
\section{Microservices}

In this section we will summarize the concept of microservices. This summary is
mostly based on \autocite{MFMS}.

Microservices are most easily explained by comparing them to, perhaps a more
traditional architecture of, monoliths. A monolithic application is typically
built as a single unit. This unit will contain all the components of which it
is built.

Increasingly the industry has been feeling frustrations towards monolithic
applications.  Changes made to a monolith usually requires the entire
application to be redeployed.  Over time it may also become increasingly hard
to keep the boundaries between components clean.  Scaling of monoliths is also
of concern, as it typically requires the entire application to be scaled, as
opposed to just the components that need additional resources.

These frustrations led to a new architectural style, eventually named
microservices. Microservices attempt to tackle each of these problems, but with
a new architecture also comes new problems.

At its core, a microservice architecture will consist of many individual and
independently deployable services. The fact that services must be designed to
be independently deployable can make it significantly easier to scale.  The
services of a system will communicate with each other through message passing.
There are no real restrictions on how services should communicate with each
other. However, in practice services will typically communicate via the
network, typically using ``dumb'' protocols, such as HTTP.

Microservices typically have a high focus on their own published interfaces.
This comes, almost automatically, from the fact that they to communicate via
message passing.

% TODO I'm not sure if we need more here.

\section{Package Managers}

In this thesis we will work solely with application-level package managers.
As mentioned in the introduction, application-level package managers deal with
application-level packages. These packages are typically software components.

The work on the Jolie Package Manager (JPM), which this thesis presents, is
hugely inspired by the following works:

\begin{enumerate}
    \item \textbf{NPM:} The package manager for Node.js\autocite{NPMC}.
    \item \textbf{Yarn:} An alternative package manager for 
    Node.js\autocite{YARNC}.
    \item \textbf{Cargo:} Build tool for the Rust language\autocite{CARB}.
    \item \textbf{Gradle:} Build tool, typically used for JVM languages
    (e.g. Java)\autocite{GRAB}.
\end{enumerate}

In the remainder of this section we will cover the most common traits from
these works.

The primary role of a package manager, as the name would suggest, is to manage
software packages. A package manager usually describes software packages
through a special file, in this thesis we will refer to this file as the
package manifest. NPM and Yarn both use a file called \txtl{package.json},
Cargo uses a \txtl{Cargo.toml}, Gradle using a \txtl{build.gradle}.
Common for all of these are, that they typically describe the package itself,
and the dependencies of this package.

Listing \ref{lst:cargo_example} shows an example for Cargo, and shows a similar
manifest for NPM is shown in Listing \ref{lst:npm_example}. Common to both of
these examples, is that their descriptions use a few attributes for describing
the package itself. Additionally they have a section describing the
dependencies, i.e., the packages that should be ``managed''. Most importantly of
the attributes describing the package are the ``name'' and ``version'', which
can uniquely refer to a specific version of the package.

\begin{listing}[H]
\begin{minted}{ini}
[package]
name = "my-cargo-crate"
version = "1.2.3"
authors = ["Dan Sebastian Thrane <dathr12@student.sdu.dk>"]

[dependencies]
log = "0.3.8"
left-pad = "1.0.1"
\end{minted}
\caption{A simple manifest for a Cargo package (a ``Crate'')}
\label{lst:cargo_example}
\end{listing}


\begin{listing}[H]
\begin{minted}{json}
{
  "name": "my-npm-package",
  "version": "1.0.0",
  "description": "This is a NPM package",
  "main": "index.js",
  "scripts": {
    "test": "echo \"Error: no test specified\" && exit 1"
  },
  "author": "Dan Sebastian Thrane <dathr12@student.sdu.dk>",
  "license": "ISC",
  "dependencies": {
    "left-pad": "^1.1.3"
  }
}
\end{minted}
\caption{A simple manifest for a NPM package}
\label{lst:npm_example}
\end{listing}

Most package managing tools come with CLI (Command-Line Interface) tools, which
help the user perform various tasks. For example, both of the manifest examples
(Listing \ref{lst:cargo_example}, and \ref{lst:npm_example}) could be generated
via their respective tools as shown in Listing \ref{lst:tools_example}.

\begin{listing}[H]
\begin{minted}{text}
dan@host:/ # cargo new my-cargo-package
    Created library `my-cargo-package` project

dan@host:/ # npm init
This utility will walk you through creating a package.json file.
It only covers the most common items and tries to guess sensible defaults.

( ... Cut for brevity ... )
\end{minted}
\caption{Package managers typically ship CLI tools which help with common
    tasks}
\label{lst:tools_example}
\end{listing}

Similarly these tools can be used to download and install their dependencies,
for Cargo \txtl{cargo build}, for NPM \txtl{npm install}. The
packages, that these dependencies resolve to, can be stored in
various places. A common approach, which we see both NPM, Yarn, and Cargo, is
to use a default centralized registry. Gradle uses a similar approach but
doesn't have a single default registry (which they refer to as repositories).
The registries will act as a database of packages, which the tools will use
to download, and install packages. Registries typically also have some type of
user management. Such a feature becomes relevant when controlling who is
allowed to publish updates for a package.

The reader may have observed already, that we have been referring to both Cargo
and Gradle as package managers, even though they were listed as build tools.
The build tool description is pulled directly from their own descriptions.
However, the border between package managers and build tools are typically
blurry. Especially because many modern tools include dependency management as a
core part of their features. Even tools, like NPM, which advertise themselves
still build provide ``build tool''-like features, such as running tests and
starting the application. In this thesis, we will simply refer to our own tool,
the Jolie Package Manager, as a package manager. This is done, since
it is its primary responsibility, and technically there is nothing to
build in Jolie, due to its dynamic nature.

\section{Introduction to Jolie}

Jolie is a service-oriented programming language and is built to support a
microservice architecture natively. In this section, we will introduce the
Jolie is, and show how it is used.

Jolie has a C-inspired syntax and is dynamically typed. Its interpreter is
written in Java.

The language has no native functions or methods, but instead, uses processes. A
process has no arguments and does not contain any stack. There are two
pre-defined processes, which will always be called by the interpreter, these
are called \txtl{init} and \txtl{main}.

\begin{listing}[H]
\begin{minted}{jolie}
include "console.iol"

define PrintOutput {
    println@Console(output)() // Prints 'OK'
}

init { a = 1 }

main {
    b = 2;
    c = a + b; // c = 3
    if (c == 3) {
        output = "OK"
    } else {
        output = "Bad"
    };
    PrintOutput // Calls the defined process 'PrintOutput'
}
\end{minted}
\caption{A very simple Jolie program}
\label{lst:simple_jolie}
\end{listing}

Listing \ref{lst:simple_jolie} shows a very simple application written in
Jolie.

The missing semicolons on lines 12, 14, or 16 are not mistakes. The semicolon
is not needed here, in fact, it would be a syntax error. The reason for this is
that the semicolon isn't used strictly for parsing purposes, but it instead
used for composing statements in a process.

Semicolon statements, or sequence statements, have a syntax of \joliel{A ; B}.
It should be read as: first, execute \joliel{A} then execute \joliel{B}. The
sequence statement requires both operands to be present, hence the syntax error
for adding a semicolon to the last statement of a block.

A similar composition operator is, the parallel statement, which has a syntax
of \joliel{A | B}, which reads as: execute \joliel{A} and \joliel{B} in
parallel.

These operators together allows the programmer to easily create a fork-join
workflow.  This is typically used in microservices when we want to collect data
in parallel, and continue once all of the data has been retrieved.

Jolie has different rules for scoping than most languages. In Jolie, everything
not defined in the global scope goes into the same scope. This also persists
through calls to defines. This is the reason that \verb!PrintOutput! can use
the output variable. To access the global scope we must explicitly state that
we wish to use it. Accessing a variable \joliel{a} in global scope would
require us to write \joliel{global.a}.

Several execution modes exists. The default execution mode, which was used in
Listing \ref{lst:simple_jolie}, is \joliel{single}. With an execution mode of
\joliel{single} the \joliel{main} process will run just a single time. Two
additional execution modes exist: \joliel{concurrent} and \joliel{sequential}.
Any of these two modes should be used if the Jolie program should act as a
service. When running in one of these two, the Jolie program will act as a
service.

When Jolie acts as a service, it will require the first statement of the main
procedure to be an input. An ``input'' represents an operation call (request)
going into the service, to which the Jolie program will need act upon.  Listing
\ref{lst:simple_input} shows a Jolie service being capable of answering
requests for either an operation called \joliel{echo} or \joliel{respondHello}.

\begin{listing}[H]
\begin{minted}{jolie}
main {
    [echo(request)(response) {
        response = request
    }]

    [respondHello(request)(response) {
        response = "hello"
    }]
}
\end{minted}

\caption{A simple service capable of answering requests for either of the two
    types of operations: \joliel{echo} and \joliel{respondHello}}

\label{lst:simple_input}
\end{listing}

Ports are the primitive that Jolie uses for communication, two types of ports
exist input and output. Ports describe a running service, where it is located
(\joliel{Location}), how to speak to it (\joliel{Protocol}), and finally which
operations it supports (\joliel{Interfaces}).

Listing \ref{lst:simple_port} shows an output port \joliel{Example}. This
output port will contact a web server running \joliel{http} on port 80 at
\joliel{example.com}. Only ports deal directly with protocols, this means that
the code interacting with the port remains the same, regardless of where it is
located, or which protocol it uses.

\begin{listing}[H]
\begin{minted}{jolie}
outputPort Example {
    Location: "socket://example.com:80"
    Protocol: http
    Interfaces: IExample
}
\end{minted}
\caption{A simple output port which contacts an example web server}
\label{lst:simple_port}
\end{listing}

The interface which the port uses is called \joliel{IExample}, a full
definition of it can be seen in Listing \ref{lst:simple_interface}. Two types
of operations exist in Jolie, namely \joliel{RequestResponse} and
\joliel{OneWay}.  The difference is fairly self-explanatory, the first receives
a request and returns a response, the other simply receives a request, and
produces no response. When an operation call is made from a client, and the
operation type is \joliel{RequestResponse}, then the client will wait for a
response to come back before continuing.  However, if the operation type is
\joliel{OneWay}, then the client will not wait for anything, and simply return
as soon as the message has been sent.

\begin{listing}[H]
\begin{minted}{jolie}
interface IExample {
    RequestResponse:
        anOperation(RequestType)(ResponseType)

    OneWay:
        hello(HelloType)
}
\end{minted}
\caption{An interface in Jolie defines which operations a port exposes}
\label{lst:simple_interface}
\end{listing}

Whenever the Jolie interpreter invokes an operation on an output port or
receives a request on an input port, the types will be checked. This check
ensures that incorrect requests aren't sent. It also ensures that we do not
attempt to process an incorrect request. Listing \ref{lst:simple_types} show
the request and response type of \verb!anOperation!.

\begin{listing}[H]
\begin{minted}{jolie}
type RequestType: void {
    .a: int
    .b: string
    .c: bool
    .d: double
    .e: any // any primitive type
}

type ResponseType: int {
    .aFixedArray[1, 3]: string
    .aNonFixedArray[0, *]: string
    .fieldWithChildren: void {
        .a: void {
            .b: int
        }
    }
}
\end{minted}
\caption{Jolie types are tree-like structures}
\label{lst:simple_types}
\end{listing}

In Jolie, types are tree-like structures, very similar to how, for example, XML
would be represented. Importantly the root may also contain a value, this is
different from how most other programming languages work. For example, a value
of type \joliel{ResponseType} is shown in Listing \ref{lst:value_example}.

\begin{listing}[H]
\begin{minted}{jolie}
value = 42; // root value
with (value) {
    .aFixedArray[0] = "A";
    .aFixedArray[1] = "B";
    .aNonFixedArray[0] = "foo";
    .fieldWithChildren.a.b = 1337
}
\end{minted}

\caption{An example of a value which would type-check for the type
    \joliel{ResponseType}}

\label{lst:value_example}
\end{listing}

Jolie natively supports a variety of techniques for composition of services.
The most important (for this work), which we will cover here are
\textbf{aggregation} and \textbf{embedding}.

Embedding allows for a larger service to run smaller services as inner
components. These services communicate with each other using more efficient
local communication. These embedded services can be other Jolie services, but
may also be services written in, for example, Java or even JavaScript.
Communicating with these services, through Jolie code, is done exactly the same
way as with any other service. As a result, it is entirely transparent to the
application code where the service is located.

Aggregation is a generalization of proxies and load balancers. An illustration
of this concept can be seen in Figure \ref{fig:aggregation}. Aggregation is
useful for creating a wide variety of proxy like architectural patterns. The
aggregation feature is often used alongside the courier and interface extender
features. Couriers allow the developer to insert code in-between the receiving
the request and forwarding it. These features, for example, allow you to add
authentication to a service which otherwise doesn't have it. This is done
entirely without having to touch the original service.

\begin{figure}[H]
    \centering
    \includegraphics[width=0.6\textwidth]{pictures/aggregation.png}
    \caption{Aggregation is a generalisation of proxies and load balancers}
    \label{fig:aggregation}
\end{figure}


\section{The Jolie Engine and Interpreter}

In this section, the internals of the Jolie engine will be introduced. This
should give the reader the necessary knowledge to understand the changes made
to support a module and package system for Jolie.

At the core of the Jolie engine is the interpreter. Each interpreter is
responsible for parsing and executing a Jolie program. A single engine may run
several instances of the interpreter, this is most commonly the case when
embedding several other Jolie services.

A simplified view of the Jolie interpreter's pipeline can be seen in Figure
\ref{fig:jolie_pipeline}.

\begin{figure}[H]
\centering
\includegraphics[width=0.8\textwidth]{background/pipeline.eps}
\caption{A simplified view of the Jolie interpreter pipeline}
\label{fig:jolie_pipeline}
\end{figure}

The first phase of the interpreter is parsing the intent. In this phase the
interpreter will determine its purpose and, if any, modifiers to its ordinary
behavior.

The Jolie engine can be invoked from the command-line, the command-line is the
source of the intent which starts the first interpreter. The syntax for the
command line is (roughly) as follows: \txtl{jolie [commands and options]
    <program> [program arguments]}.

Options change the overall behavior of the engine, all interpreter instances
share these. These change the behavior of the interpreter.

Commands and program arguments, only belong to the interpreter that they were
originally passed. These determine what the interpreter will do.

The intent parsing phase is responsible for locating and retrieving the program
used.

A URI matching the input program is given to the program parser, the second
phase. The parser is responsible for creating the Abstract Syntax Tree (AST)
    which represents the input program. The parser will produce only a single
    root node, namely the \txtl{Program} node. Includes are thus implemented by
    essentially inserting the contents of the include target. A quirk for Jolie
    includes, is that relative file includes are not resolved relative to the
    current file, but rather to the current working directory, i.e., where the
    engine was started.

The third phase traverses the AST to make sure that it is semantically valid.
This weeds out programs that are syntactically correct, but not semantically
valid. The checks performed in this phase are somewhat limited, given the
dynamic nature of the language.

Given a semantically valid AST the interpreter is ready to build the
interpretation tree. The interpretation tree contains new nodes, which are
runnable.

Once the interpretation tree is built, we're ready to execute the actual code
(which lives in the interpretation tree). The \txtl{CommCore} component is
responsible for all communication, we'll discuss this briefly when relevant.
Following its initialization, the user code of Jolie will start executing. The
\joliel{init} procedure is executed first followed by the \joliel{main}
procedure.

\section{A Complete Application with Jolie}
\label{sec:complete}

In this section, we will describe a complete application written in Jolie. The
application will contain several services and will be written using best
practices from before the module and package system.

Along the way, several patterns used for re-using code will be highlighted.
This will serve as the core motivation for our solution. The observations made
in this section are drawn from the Jolie Documentation\autocite{JOLB}, the
Leonardo\footnote{A web server written in Jolie by the Jolie team. Used to
    power the official Jolie website} web server\autocite{LEO}, and the
    source-code for the Jolie website\autocite{LEOB}.

\subsection{Architecture}

% Things we wish to cover in this section:
%
%   - Includes and why they will be problematic
%   - Configuring location of ports, both external and embedding of stuff
%   - Project structure, and why this is a problem
%   - Parameter configuration

The application we'll be building is a very simple calculator system. The
architecture of our application is shown in Figure \ref{fig:complete_arch}.

The calculator service will provide various operations for applying an operator
on a sequence of numbers. In order to perform this, the action of applying
these operators, the calculator will contact other microservices. This example
is obviously contrived. Writing services for such simplistic services in the
real world would be considered an anti-pattern. However, this example allows us
to see services interact with each other without thinking that much about the
business logic of the services, which in this case is fairly irrelevant.

\begin{figure}[H]
\centering
\includegraphics[width=1.0\textwidth]{background/calc_system.eps}
\caption{The Architecture of a Simple Microservice System}
\label{fig:complete_arch}
\end{figure}

\subsection{Implementing the Calculator Service}

We will keep our focus on the \txtl{calculator} service, and its interactions
with other services. Putting together the stuff learned from the previous
sections, we can quickly create an input port for the service with the
appropriate interface.

It is considered best practice to place the public interfaces that a service
exposes in its own separate file. Files intended for other services to include
typically have the file extension \txtl{.iol} as opposed to \txtl{.ol}. There
is no technical difference between the two, but it allows for the developer to
more easily express intent. Files that are included can be placed in the
\txtl{include} directory. This directory is always implicitly added to the
search path.

Thus in order to implement our calculator service we create two files. One file
for the service implementation (\verb!calculator.ol!), and another which can be
used by other services (\verb!calculator.iol!). These files are shown in
Listing \ref{lst:simple_start}.

\begin{listing}[H]
\begin{multicols}{2}

\begin{minted}[breaklines,fontsize=\footnotesize]{jolie}
// calculator.ol
include "console.iol"

inputPort Calculator {
    Location: "socket://localhost:12345"
    Protocol: sodep
    Interfaces: ICalculator
}

main {
    [sum(request)(response) {
        println@Console("Implementation goes here")()
    }]

    [product(request)(response) {
        println@Console("Implementation goes here")()
    }]
}
\end{minted}

\columnbreak

\begin{minted}[breaklines,fontsize=\footnotesize]{jolie}
// calculator.iol
type Numbers: void {
    .numbers[2, *]: int
}

interface ICalculator {
    RequestResponse:
        sum(Numbers)(int),
        product(Numbers)(int)
}
\end{minted}
\end{multicols}
\caption{Interface and implementation of the \txtl{calculator} service}
\label{lst:simple_start}
\end{listing}

\begin{observation}

The interfaces and types of a service are typically separated into their own
files. These files typically have the extension \verb!.iol! to indicate that it
does not contain a service implementation, but is rather intended for inclusion
by a service which requires it.

These files are typically put in the \verb!include! directory. This
directory is always added to the search path natively by the Jolie engine.

\end{observation}

The operations which the calculator exposes needs to collaborate with the other
services. In order for us to speak to them, they need output ports.  The output
ports need interfaces, just like the calculator service these should be exposed
from a dedicated file.

Secondly, these output port needs to be reached. The port can either be
embedded or we can provide external bindings to it. For output ports, we may
change this binding dynamically at runtime. This is unlike input ports which
must be ready at deployment time.

Binding an output port to an external service, is done by setting the
\joliel{Location} and \joliel{Protocol} properties. For example, to make
\joliel{Multiplication} bind to an external service we could define it as shown
in Listing \ref{lst:mult_ext_port}.

\begin{listing}[H]
\begin{minted}{jolie}
outputPort Multiplication {
    Location: "socket://mult.example.com:443"
    Protocol: https
}
\end{minted}
\caption{The \joliel{Multiplication} port is bound to an external service}
\label{lst:mult_ext_port}
\end{listing}

It is possible to define embeddings of services in the \joliel{embedded} block.
For our desired deployment, we wanted to embed the addition service inside of
the calculator service. Assuming that the addition service is written as a
Jolie service, and its service is implemented in \txtl{addition.ol}, then we
may create an embedding as shown in Listing \ref{lst:simple_embedding}. Just
like in the case of the external services, the input port of the receiving
service \emph{must match}. In the case of embedded service, there must be an
input port listening on the \joliel{"local"} location.

\begin{listing}[H]
\begin{minted}{jolie}
embedded {
    Jolie:
        "addition.ol" in Addition
}
\end{minted}

\caption{Embedding the \mintinline{jolie}{addition} service in the
    \mintinline{jolie}{Addition} output port}

\label{lst:simple_embedding}

\end{listing}

Quite often the location of an input port is considered a deployment problem.
We see this quite clearly in the case of embedding a service. All solutions to
change the deployment of a service in Jolie, all require us to \emph{modify the
    source code of a service}.  Best practices in Jolie attempt to make this
    less of a problem by including a configuration file which contains
    constants. The \txtl{addition} service, might include a file called
    \verb!addition.iol! which sets up the location and protocol of the service.
    An example of this is shown in Listing \ref{lst:include_as_conf}.

\begin{listing}[H]
\begin{multicols}{2}

\begin{minted}[breaklines,fontsize=\footnotesize]{jolie}
// addition.ol
include "addition.iol"

inputPort Addition {
    Location: ADDIITON_LOCATION
    Protocol: ADDIITON_PROTOCOL
}
\end{minted}

\columnbreak

\begin{minted}[breaklines,fontsize=\footnotesize]{jolie}
// addition.iol
constants {
    ADDIITON_LOCATION = "local"
    ADDIITON_PROTOCOL = sodep
}
\end{minted}

\end{multicols}

\caption{A common Jolie practice for solving configuration of a service, is to
    include a file containing constants with the desired configuration.}

\label{lst:include_as_conf}

\end{listing}

\begin{observation}

Most configuration of Jolie services is done via constants. Jolie constants,
different from most other languages, can be simple literal types, or they
may even contain identifiers. This makes them rather versatile in what
they can configure.

The constants are typically placed in a separate source file, which is
included by the service requiring the configuration.

\end{observation}

In Listing \ref{lst:simple_start} the output port \txtl{Console} is used. This
output port isn't visibly defined in the source code we see, it is instead
defined in the included file \txtl{console.iol}. This is common practice in
Jolie. This is especially used for services that work in a library-like
fashion, that is not intended as a stand-alone service. This pattern is used
for almost all services in the Jolie standard library\autocite{JOLB}.

\begin{observation}

Services intended to be used as libraries are often contained in a single
\verb!.iol! file. This file contains everything required to set it up, included
interfaces, types and an embedded output port.

This makes it very easy to use the service, however, it also makes it
impossible to bind to this service externally, without first getting an
embedding. This isn't possible since we cannot include the interfaces and
types by themselves.

\end{observation}

With the output ports correctly configured, we may now implement the actual
business logic for our calculator. This might look like shown in Listing
\ref{lst:op_impl}.

\begin{listing}[H]
\begin{minted}{jolie}
[sum(request)(response) {
    total = 0;
    for (i = 0, i < #request.numbers, i++) {
        // Add current number with total, and store result in total
        add@Addition({
            .a = total,
            .b =  request.numbers[i]
        })(total)
    };
    response = total
}]
\end{minted}
\caption{Implementing the checkout operation}
\label{lst:op_impl}
\end{listing}

Finally, we want to reflect slightly on the file structure that this service
has. In order to use external services, we had to include files which contain
these interfaces. This version is worse when an embedding is desired, since the
entire service with its implementation is suddenly required.

The files from these external services are also hard to manage. It isn't
possible to simply move the source code of a service into its own directory.
This is not possible since includes, unlike in most other languages, are not
relative to the file performing the include, but rather relative to the project
level root. Thus if we wish to move the service to a new directory, all
includes in the source code would have to be updated. A common result of this
is that every single file gets dumped into the project level root. When all
files are stored in the same directory, we also get a much larger chance of
having name conflicts. As a result, names tend to become rather large, to make
it less likely that a conflict occurs.

The final file structure of the calculator service is shown in Listing
\ref{lst:file_structure}.

\begin{listing}[H]
\begin{minted}{text}
.
|-- addition.iol
|-- addition.ol
|-- calculator.ol
|-- include
|   `-- calculator.iol
`-- multiplication.iol
\end{minted}
\caption{File structure of the calculator service}
\label{lst:file_structure}
\end{listing}

\begin{observation}

    In order to collaborate with other services, it is often needed to manually
    copy files into the service. Which files are required depend on the
    deployment. When an embedding is required, we will need the entire service
    (and any services it may depend on itself). If we just need to interface
    with it, we can simply include the files required for its interface files.

    Structuring the files is problematic, due to \verb!include! statements not
    being relative to file performing inclusion.

\end{observation}




\chapter{Jolie Modules}
\section{Summarizing the Problem}

The goal of this thesis is to promote, and facilitate better code reuse in
Jolie. In Section \ref{sec:complete} we showed how a typical Jolie service
would be written, and summarized a number of patterns we have observed.


\section{Modules: Design and Implementation}

The module system of Jolie is made to facilitate better options for the
modularization of code. This is what allows for a package-manager to exist. We
discuss the package manager in detail in Chapter \ref{cha:packages} and
\ref{cha:pkg_mngr}. A Jolie module is described by the following three
properties:

\begin{enumerate}

    \item \textbf{A name.} For a single running Jolie engine, this will
    uniquely identify a module

    \item \textbf{A file path to the Jolie code-base.} We will refer to this as
    the project root or just root.

    \item \textbf{(Optional) The entry point to the Jolie code.} This is a file
    path given relatively to the project root.

\end{enumerate}

The Jolie engine needs to know about these modules. The engine is informed
about these modules from the intent, passed as options, which starts the
engine. This information is collected in the ``intent parsing'' phase and is
made available to any later phase that might need it.  Since the intent is
passed as an option, all interpreter instances inside the engine will know
about each module.  As a result, any embedded service will also be aware of the
same modules.

Consider a module ``foo'', with its source code placed at \txtl{/packages/foo}
and an entry-point at \txtl{/packages/foo/main.ol}. It is possible to inform a
Jolie interpreter about this module using
\txtl{--mod foo,/packages/foo,main.ol}.

As a result, a lot of additional options must be passed to the engine. In this
case, a package manager is expected to take the heavy lifting and figure out
which modules exist.  This allows for more freedom in how these tools are
implemented. If another approach were to be found later, then a new
implementation could simply be created without breaking the language.

A new include primitive has been added, a module include. This has been created
to allow for better internal organization. This primitive allows the developer
to include files from a particular module. The new primitive is shown in
Listing \ref{lst:mod_include}. When a module include is used, the search path
will be altered, such that the project root is changed to that of the package.
As opposed to the original include primitive, which would look for files
relative to the current working directory.

The new include primitive, along with native support for modules, allows for
the code to be properly organized. With this, a module may be placed inside of
its own directory, without any code has to change.

\begin{listing}[H]
\begin{minted}{jolie}
include "<file>" from "<module>"
\end{minted}
\caption{Extension to the include statement, made for module imports}
\label{lst:mod_include}
\end{listing}

\subsection{Module Include Algorithm}

A module include is an extension of ordinary includes which specify which
module to perform the inclusion from, that is \joliel{include "file" from
"module"}.

The base case of such an include is simple. We make a ``context switch'' from
our working directory into the root of the module and start searching for the
file in that folder. This works quite well, however, consider the following
scenario.  From a client we perform a module include: \joliel{include
    "interface.iol" from "dependency"}.

\begin{listing}[H]
\begin{minted}{java}
/* /dependency/interface.iol */
include "types.iol"

interface Dependency { /* ... */ }

/* /dependency/types.iol */
type SomeRequest: void { /* ... */ }
\end{minted}

\caption{Module includes may cause more ordinary includes. We must make sure
    the switch in context is kept}

\label{lst:dependency_include}
\end{listing}

With the code in Listing \ref{lst:dependency_include}, the system would
incorrectly look for \txtl{types.iol} in the current working directory! To fix
this a stack is kept.  This stack gains an entry when a module include is
performed, this entry is popped off the stack once the contents of that include
has finished parsing.  The stack is then used to inject the module's root into
the search path for the ordinary includes.

\section{Configuration}

% TODO Define configuration and which constructs we want to be configurable

\subsection{Motivation}

From observation 2 we saw that most configuration was done via the inclusion of
source code. This source code would expose constants (read: literal values
\emph{and} identifiers). The included source code, however, can do anything
that Jolie source normally can, and isn't limited to just the desired
configuration. As a result, a service developer cannot be certain that the
configurator (entity who provides configuration) doesn't start messing with
other details of the program. Deploying defensive
programming\footnote{Defensive programming techniques are usually employed for
systems that require high availability, or where safety and security is
required.} techniques against this becomes significantly more problematic,
since no guarantees about the configuration source file can really be made.

Distributing re-useable packages is also problematic with this approach.
Several features of package management requires that the package remains
read-only. For example, updating of packages require this, without it
source-code merges would be required.

We also saw that a lot of issues, that should have been a deployment issue,
became a code issue.

This gives us plenty of reason to explore the need for a native configuration
format.  Most other systems would most likely go for a system defined in user
code, as opposed to natively. An example of such framework, could be Vert.x, it
is a tool-kit for building reactive applications on the JVM.  Examples of such
``reactive applications'' are microservices. The configuration workflow is
shown in Figure \ref{fig:normal_conf}. The system will retrieve, and read
external configuration files, directed by the user code, and apply the
configuration as needed.

% Most other systems can do this at run time
% For example Vert.x does this by reading external configuration files
% Once configuration is done, we can start up the server

\begin{listing}[H]
\begin{minted}{text}
+------------------+     +-------------------+     +--------------------+
| executable start | --> | retrieve conf     | --> | read conf files    |
+------------------+     +-------------------+     +--------------------+
                                                             |
                            +-----+                          |
                            |     | reconfigure              |
                            |     v                          v
                         +-------------------+     +--------------------+
                         | server running    | <-- | perform conf       |
                         +-------------------+     +--------------------+
\end{minted}
\caption{Simplified workflow for configuration of Vert.x applications}
\label{fig:normal_conf}
\end{listing}

% http://vertx.io/blog/vert-x-application-configuration/
% http://vertx.io/docs/vertx-config/kotlin/

TODO Defined deployment time somewhere

However implementing such as a system in Jolie has its problems, most of these
come from the difference between general-purpose programming languages and
specialized programming languages.

In general-purpose languages, the constructs (such as a server's socket) for
the microservice architecture are created in user code. As a result they are
entirely accessible from user code. This make it feasible to change their
behaviour, since code can run before deployment occurs.

In Jolie the constructs are managed directly by Jolie. Doing this has multiple
advantages, such as less complexity in user code, but it also means that user
code is capable of doing less. Jolie user code can for example not control
networking directly, but is instead forced to use the abstractions provided by
Jolie (sending messages). The language puts constraints on certain
constructs being fully prepared directly in the source code. This is analogous
to a programming language requiring the types of a struct's field to be present
at compile time. As a result, not all constructs can be changed at run time.
Concrete examples of this includes the input ports, which needs to be ready at
deployment time. Thus without native support for configuration of these, it
would not be possible to change the input port.

\subsection{Configuration Units}
\label{sec:conf_units}

A configuration unit is the basic entity, which encapsulates the configuration
of a single Jolie module. A configuration unit is known by its name (known as
its ``profile''), along with which module it configures. Having multiple profiles
for the same module can be useful for a variety of use-cases. A common
use-case, could for example be to have separate profiles for development and
production.

The units hold configuration for every possible type of configurable construct
in Jolie. The ones supported are:

\begin{enumerate}
    \item Input and output ports
        \begin{itemize}
            \item Location
            \item Protocol and protocol parameters
            \item Embedding of other services (output ports only)
        \end{itemize}
    \item General purpose parameters
    \item Interface rebinding
\end{enumerate}

In the coming sections we'll mostly focus on the first two, in Section TODO
we'll cover interface rebinding.

% TODO Do we allow for inclusion of other files?
Configuration units are defined in configuration files, which may contain
several units. These files may even include other files, to pull in more
configuration units.

Jolie provides a new configuration file format. This format is custom, and made
to mimic the syntax of Jolie. Listing \ref{lst:simple_conf} shows a very
simple configuration unit. This units sets the location and protocol for the
output port \mintinline{jolie}{A}, the location of the input port
\mintinline{jolie}{ModuleInput}, and a parameter.

\begin{listing}[H]
\begin{minted}{jolie}
profile "hello-world" configures "my-module" {
    outputPort A {
        Location: "socket://a.example.com:3000"
        Protocol: sodep { .keepAlive = true }
    },

    inputPort ModuleInput {
        Location: "socket://localhost:80"
    },

    myParameter = 42,
    myParameter.subProperty = "hello"
}
\end{minted}

\caption{A simple configuration unit named \mintinline{jolie}{hello-world}
    configuring the module \mintinline{jolie}{my-module}}

\label{lst:simple_conf}

\end{listing}

Embedding of output ports can be performed from within a configuration unit.
This moves the embedding from being a code problem to, what it should have
been, a deployment problem. Listing \ref{lst:conf_embedding} shows the
embedding of output port \mintinline{jolie}{A}. Note that we need to make a
reference to the module, since the profile names are placed under a namespace
for each module. This way multiple services can share the same name, a
situation which is likely to occur with common profile names, such as
``development'' and ``production''.

\begin{listing}[H]
\begin{minted}{jolie}
profile "hello-world" configures "my-module" {
    outputPort A embeds "a-module" with "a-profile"
}

profile "a-profile" configures "a-module" {
    // configuration of a-module goes here.
}
\end{minted}
\caption{Embeddings make reference to other configuration units}
\label{lst:conf_embedding}
\end{listing}

As we can see from the examples, it is not necessary to provide all the values
of a port. It isn't necessary for two reasons. The first reason is that certain
values may be provided by the underlying module, which uses this unit. If a
module provides a value, then the configuration unit cannot override it. The
second reason is that configuration profiles may extend other profiles.

Configuration units may extend another unit, which configures the same module.
The tree of inheritance may be of an arbitrary depth, but each unit may only
extend a single unit, and they must configure the same module. The child is
also wins when it comes to configuration. This means that if unit ``B'' extends
``A'', and they both configure the same value, then the values found in B is
the one that is correct. Listing \ref{lst:conf_extends} shows an example of
extension with units.

\begin{listing}[H]
\begin{minted}{jolie}
profile "a" configures "a-module" {
    aValue = 42,
    aValue.sub = "hello",

    outputPort ExternalService {
        Location: "socket://external.example.com:42000"
    }
}

profile "b" configures "a-module" extends "a" {
    aValue = 100
    // aValue.sub = "hello"
    // ExternalService.location = "socket://external.example.com:42000"
}
\end{minted}
\caption{Configuration units may extend other units}
\label{lst:conf_extends}
\end{listing}

The module developer is often aware of what the defaults should be. For this
reason default configuration profiles may be shipped along the modules, which
are implicitly imported into every configuration file. The Jolie engine will
look for any \mintinline{text}{.col} file\footnote{The file extension of the
configuration file} in the \mintinline{text}{conf} folder. This folder
should be placed relative to the module's root. For example, if module "a" has
a file called \mintinline{text}{conf/my-defaults.col}, which contains a unit
called "default". Then the user of the package may either write a configuration
unit which extends this, simply by writing \mintinline{jolie}{profile
"something" configures "a" extends "default"}, or the default
directly. There is no need for any inclusion of this file.

It should be noted that no single unit is required to provide all
configuration. The system doesn't have any ``abstract''\footnote{As in abstract
classes, a concept often used in object oriented programming} configuration
units. However it is required that the configuration file provides all the
necessary configuration, as declared by the module. We'll learn more about
how a module declares configuration in Section \ref{sec:ol_conf}.

\subsection{Configuration and the Core Language}
\label{sec:ol_conf}

In section \ref{sec:conf_units} we learned about the configuration files, but
never actually saw the corresponding Jolie modules. In this section we'll
discover how the language has changed to accommodate this configuration
feature.

The configuration system provides the ability to configure any type of port,
and additional provide the module with ``parameters''. Ports were covered in
section TODO. Most importantly it should be noted that output ports do not
require any values at deployment time, and can be changed at run time, while
input ports require all their values at deployment time, and cannot be changed
at run time. % TODO Do we want to cover why it is hard to change input ports at
% run time?

\begin{listing}[H]
\begin{multicols}{2}
\begin{minted}[breaklines,fontsize=\footnotesize]{jolie}
outputPort Service {
    Location: "socket://service.example.com:443"
    Protocol: https {
        .debug = true,
        .debug.showContent = false
    }
    Interfaces: IService
}
\end{minted}

\columnbreak

\begin{minted}[breaklines,fontsize=\footnotesize]{jolie}
outputPort Service {
    Interfaces: IService
}
\end{minted}

\end{multicols}

\caption{
    Two valid output ports.
    (Left) A fully configured output port.
    (Right) A minimal output port.
}

\label{lst:basic_output}
\end{listing}

Listing \ref{lst:basic_output} shows a two different output ports, both
perfectly valid.

TODO

When should we configure? Go through a few different approaches, always
override source, always let code win, do we let everything be configurable?
Should configuration stay final, or should we allow changing it later?  How do
we distinguish between accidentally leaving something out and wanting it to be
configurable?

Output ports

Input ports, and why they are harder.

The \mintinline{jolie}{dynamic} keyword has been introduced to deal with which
ports that should allow configuration, any port which does not have the dynamic
keyword is static. Only ports which are static should be configurable. TODO Why
is this a good idea.

Parameters

\subsection{Implementing the Configuration Format}

Implementation, and fixing the AST

\subsection{Example}

TODO This is just copy pasted from some markdown document. Might be able to
find a better example. Syntax is most likely also outdated.

\includegraphics[width=\textwidth]{prototypes.png}

(The dashed region corresponds to which services that are embedded together)

This use-case is intended to show how we can easily go from a prototype where
we embed everything (this is easier to run locally) to hosting each service by
itself.

\begin{minted}{jolie}
// A.ol

ext outputPort A {
    Interfaces: AIface
}

ext outputPutPort B {
    Interfaces: BIface
}

constants {
    FOO: int
}
\end{minted}

\begin{minted}{jolie}
// A.col

include "B.col" // The includes work just like they do in Jolie
include "C.col"

// Previously called namespace. Configures makes it more explicit that we're
// talking about a specific package and not an arbitrary name
configures "A" {
    // Like always we just put the definitions here
    FOO = 42

    // We alter the syntax slightly for output ports being embedded
    outputPort B embeds B
    // The second B refers to the profile B (if no profile is specified it gets
    // the same name as the package it configures). This means that this
    // configures block also has name "A". It could also have been written as:
    // `profile "A" configures "A"`

    outputPort C embeds C
}
\end{minted}

Note: The previous proposal did not allow for embedding of services directly
from the configuration. This is however needed since packages by themselves
should be considered read-only. Thus if we want to configure a package to embed
its dependencies, then this must be done from external configuration, we cannot
do this in source.

\begin{minted}{jolie}
// B.col

configures "B" {
    inputPort B {
        Location: "local"
    }
}
\end{minted}

\begin{minted}{jolie}
// C.col

configures "C" {
    inputPort C {
        Location: "local"
    }
}
\end{minted}

In order to use an external B we need to update "A.col":

\begin{minted}{jolie}
// A.col
include "C.col"

configures "A" {
    outputPort B {
        Location: "socket://b.example.org:8000"
        Protocol: sodep
    }
    outputPort C embeds C
}
\end{minted}

In order to update the deployment file of B we simply need to update the
input port to no longer be local.

\begin{minted}{jolie}
// B.col
configures "B" {
    inputPort B {
        Location: "socket://localhost:8000"
        Protocol: sodep
    }
}
\end{minted}

In most cases however this would be unnecessary since the default
configuration file for "B" could already include a default input port.

In that case we could simply deploy directly from the default configuration.
The restriction on configuration units to only configure a single package helps
a lot. This restriction means that we cannot from any node configure any other
node which isn't a direct child of it. Without this we wouldn't be able to
easily swap out one configuration unit for another.


\section{Implementing the Configuration Format}
\label{sec:col_impl}

The new interpreter pipeline is as follows, new phases marked in bold text:

\begin{enumerate}
\item Intent parsing
\item Jolie code parsing
\item \textbf{Configurator}
\item Semantic verifier
\item Interpretation tree builder
\item \textbf{Late verification}
\item Communication core init
\item Init \& main process
\end{enumerate}

The primary work occurs in the configurator service, but some of the
verification work, namely the parameters, are postponed until after the
interpretation tree has been build.

\subsection{The Configuration Phase}

The configuration phase builds upon the result of the intent and parsing phase.
From the intent phase, we receive information about modules and the selected
configuration unit. From the parsing phase we, obviously, get the AST.

The configuration phase works by modifying the AST to insert the configuration
gathered from the selected configuration unit. The processed AST is handed to
the later phases, and the remainder of the pipeline works mostly like it has
before.

\subsubsection*{Processing of Configuration Files}

\paragraph{Resolving configuration files}

The first step of the configurator is to gather a fully processed configuration
unit. This process starts by looking for the configuration file which was
passed from the intent. The selected configuration unit is passed with:
\\\txtl{--conf <unitName> <filePath>}.

If the file path is an absolute file path, then that file will always be used.
However if the file path is relative some searching is required. The base case
for resolving this is to look for the configuration file relative to the
current working directory. However this approach stops working once modules starts
embedding other modules in a static fashion.

Consider three modules \txtl{A} and \txtl{B}, \txtl{C}. Module \txtl{A} depends
on \txtl{B}, and \txtl{B} depends on \txtl{C}.

\begin{listing}[H]
\begin{minted}{jolie}
/* /a/main.ol */
outputPort B { Interfaces: BInterface }

/* /a/config.col */
profile "a-deployment" configures "a" {
    outputPort B embeds "b" with "b-deployment"
}

profile "b-deployment" configures "b" { /* ... */ }

/* /b/main.ol */
outputPort C { Interfaces: CInterface }

embedded {
    Jolie:
        "--conf c-deployment embedded.col c.mod"
}

/* /b/embedded.col */
profile "c-deployment" configures "c" { /* ... */ }
\end{minted}

\caption{Relevant source code for modules \txtl{A} and \txtl{B}}

\label{lst:need_for_parent}

\end{listing}

Relevant code is shown in Listing \ref{lst:need_for_parent}.  Assume that we
start module \txtl{A} with \\\txtl{--conf a-deployment config.col a.mod}.  The
problem arises because the fixed embedding of module \txtl{C} performed in
\txtl{B}. Because the files are resolved relative to the working directory in
which the engine was started, the engine would look for the configuration file
in module \txtl{A}. To fix this, every time an embedding is performed the
engine will implicitly add a \txtl{--parent} flag. This flag informs the new
interpreter who performed the embedding. If the parent flag has been given, the
configuration file will be found relative to that of the parent.

The working directory, throughout this, is \txtl{/a/}.  The following will
happen:

\begin{enumerate}

\item Start engine (from CLI) with: \txtl{--conf a-deployment config.col a.mod}
File resolved to \txtl{/a/config.col} because of relative path and no parent.

\item Configuration causes embedding of \txtl{b-deployment} with:
\\\txtl{--parent a --conf b-deployment /a/config.col b.mod}. File resolved to
\\\txtl{/a/config.col} because of absolute path.

\item Source embedding in \txtl{B} causes embedding of \txtl{c-deployment}
with: \\\txtl{--parent b --conf c-deployment embedded.col c.mod}. File resolved
to \\\txtl{/b/embedded.col} because of relative path with parent.

\end{enumerate}

\paragraph{Parsing of configuration files}

The resolved configuration file is passed to the configuration parser along
with known modules. The parser is written in the same style as the Jolie code
parser.  The parser is a hand-written recursive descent parser with one token
lookahead.

The parser directly outputs the AST. At the top level, the AST consists of
\javal{ConfigUnit}s which directly map to the configuration unit abstraction.
Each \javal{ConfigUnit} is namespaced under the module that they configure.
Internally each namespace maps the name of the \javal{ConfigUnit} to the actual
\javal{ConfigUnit}. The top-level node of the AST contains a dictionary mapping
between the name of modules and their namespaces. The Java type for this
mapping ends up being \javal{Map<String, Map<String, ConfigUnit>>}, this allows
for easy access of \javal{ConfigUnit}s when they are needed. This does however
also mean that duplicates must be caught during parsing since they would
otherwise just silently override previous entries. This would typically have
been performed in a later stage to keep the parser ``pure''. In this case, some
purity was sacrificed for some efficiency and ease of use.

The grammar of the configuration format closely follows the grammar of Jolie
code. As a result, the configuration format also re-uses AST nodes where
possible. As we will see later, this proved helpful when outputting a
configured AST. The configuration grammar only supports a subset of the
features that Jolie code does.

Recall from Section \ref{sec:conf_units} that each module may publish default
configuration units, which can be used by user-level configuration units. For
this reason, the parser will make special note of references to other
configuration units. Such references currently only appear in output port
embeddings, e.g., \\\joliel{outputPort Foo embeds "foo" with "a-foo-unit"}.

Once the initial file has been parsed, every single configuration file placed
in the ``conf'' directory of referenced modules are parsed and placed in the
same AST. Configuration files parsed during this phase may itself have
references to other modules, in that case these will be pulled in as well.  A
set of parsed defaults is kept to avoid parsing the same file multiple times
or potentially reach infinite loops.

\paragraph{Merging configuration files}

With the entire configuration tree in place, merging is performed to create a
single \javal{ConfigUnit} which represents the actual configuration of this
module. The merging occurs between a \txtl{child} and its \txtl{parent} unit
(that is, the one it extends, if any) and results in a \txtl{merged} unit. The
\txtl{parent} is always merged, via a recursive call, first. The first
\txtl{child} selected is the \txtl{ConfigUnit} matching the profile name and
package given by the intent.

For ports, the merging algorithm works by inserting all the ports of the
\txtl{child} into the \txtl{merged} unit. The ports of the \txtl{parent} are
then merged with the preliminary output of the \txtl{merged} unit. The
process of merging a \txtl{port} of the \txtl{parent} with the \txtl{merged}
unit is shown, in pseudo code, in Listing \ref{lst:merge_ports}. In short the
merging will favor the \txtl{child}, and the \txtl{parent} will only provide
values if the \txtl{child} doesn't.

\begin{listing}[H]
\begin{minted}[linenos=true]{text}
procedure merge(port: Port, merged: ConfigTree): Port {
    if (!merged.contains(port)) return port
    child = merged.getExisting(port)

    if (child.isComplete) return child
    // The child cannot be an embedding, since those cannot
    // be partial.

    for (property in Port.properties) {
        if (child[property] == null) {
            child[property] = merged[property]
        }
    }
    return child
}
\end{minted}

\caption{Pseudo code for merging a \txtl{parent}'s \txtl{port} into a
    \txtl{merged} configuration tree}

\label{lst:merge_ports}

\end{listing}

For parameters the same idea applies. However each parameter definition node in
the AST potentially only provides a partial value. For example the node
\joliel{foo.bar = 42}, where the type of \joliel{foo} is \joliel{string { .bar:
    int }}, only provides a partial definition of \joliel{foo}. For this
    reason, the parameter assignments are first added from the \txtl{parent}
    and then by the \txtl{child}. This allows the definitions of the
    \txtl{child} to override those of the \txtl{parent}.

\subsubsection*{Applying Configuration to the AST}

The \txtl{merged} unit from the previous section is then used along with the
AST corresponding to the Jolie program. The last phase of the configurator will
process the entire AST and output a new, fully configured, AST. This process
works by running a \txtl{process} function on every single top-level node. The
\txtl{process} function will inspect the node and return a list of replacement
nodes.

Throughout this process the configurator will also verify that no configuration
is performed on nodes which aren't configurable. This includes dynamic ports
and non-empty interface definitions.

\paragraph{Ports}

When a port node is encountered in the AST, the \txtl{merged} unit is searched
for an output port matching its name. If none is found, the node is returned as
it is. Otherwise the configuration proceeds. For input ports and non-embedding
output ports, the process is the same. The properties of the AST port and the
port represented in the configuration unit are compared one-by-one. If any
property is defined by both, an error is returned pointing to the AST node and
the place that the configuration took place. At the end, if no errors are
returned, a new port is returned with their properties merged together.

For output ports which embed another module, the process is slightly different.
The configurator will first verify the existence of the configuration unit by
looking it up in the configuration tree. If it is not found an error will be
returned. Otherwise, the configurator will output a new embedding node. The
embedding node will contain a reference back to the same configuration file the
current \txtl{ConfigUnit} comes from. When the embedding is performed the
configurator of the embedded service will lookup the correct configuration unit
from the same file. This will cause an addition parsing of the same file, but a
caching layer at the configuration parser could take care of this.  The
embedding process is demonstrated in Listing \ref{lst:embedding_output}.

\begin{listing}[H]
\begin{multicols}{2}

\begin{minted}[breaklines,fontsize=\footnotesize]{jolie}
embedded {
    Jolie:
        "--conf foo-profile <inputConfigFile> foo.mod" in Foo
}
\end{minted}

\columnbreak

\begin{minted}[breaklines,fontsize=\footnotesize]{jolie}
outputPort Foo embeds "foo" with "foo-profile"
\end{minted}

\end{multicols}

\caption{AST nodes generated, shown as code (left), by configuration (right)
    for an output port containing an embedding}

\label{lst:embedding_output}

\end{listing}

\paragraph{Interface Rebinding}

Just like ports, when an interface definition is found, the \txtl{merged} unit
is searched for a matching interface definition. Assuming no errors, the
configurator must now replace the dummy interface definition with the real
interface definition, as the configuration specifies.

Quickly recapping interface rebinding in the configuration is shown in Listing
\ref{lst:interface_rebinding_recap}. Note that the only information the
configurator has for locating the replacement is the name of the interface and
its module.

\begin{listing}[H]
\begin{minted}{jolie}
interface Dummy = Concrete from "my-module"
\end{minted}

\caption{Rebinding the interface \txtl{Dummy} to match the interface
    \txtl{Concrete} from the module \txtl{my-module}}

\label{lst:interface_rebinding_recap}

\end{listing}

The interface will be found by parsing the module, starting at its entry point.
Its AST, if it exists, will then be searched for an interface matching the
replacement. Simply replacing the dummy with its replacement, however, isn't
enough. Consider, for example, the scenario shown in Listing
\ref{lst:interface_rebinding_types}. In this scenario, the processed AST will
no longer be valid due to the missing types (\txtl{FooRequest} and
\txtl{FooResponse}).

\begin{listing}[H]
\begin{minted}{jolie}
type FooRequest: void {
    .child: string
}

type FooResponse: int

interface Concrete {
    RequestResponse:
        foo(FooRequest)(FooResponse)
}
\end{minted}

\caption{Simply copying the interface definition is not enough, the types must
    also be copied}

\label{lst:interface_rebinding_types}
\end{listing}

The relevant type definitions are found by iterating through every operation of
the replacement interface. If the types listed in the operations are type links
those are followed. This leads to a recursive algorithm, which looks through
every type looking for type links, which will cause more types to be copied
over. In order to avoid infinite loops a set of already visited types is kept.
If a type has been seen before then its children will not be visited.

In the end, the original interface definition will be replaced, and all
necessary type definitions will also be added to the program.

\paragraph{Parameters}

Parameters being a new concept in Jolie simply outputs new nodes for every
single parameter assignment. This along with the definitions gathered from
ordinary Jolie code parsing are processed in the interpretation tree building
phase.

\subsection{Additions to the Semantic Verifier}

As we have seen, several syntax changes were also made to the core languages.
These modifications to the syntax changed which programs were
legal. As an example, interfaces are now allowed to be empty at a syntactical
level, to signal that they should be replaced by configuration. However having
an empty interface not being replaced should still cause an error.

This means that additional checking must now be performed, since it no longer
will be caught during parsing, to ensure that programs are semantically valid.
Such a phase already exists in the compiler, the semantic verifier.

The semantic verifier works by performing an AST traversal and validate the
semantics of each node.

The verifier has been updated to not allow empty interfaces, this was
previously enforced at a syntax level. Since the configurator runs before this
phase and updates the AST, this will only cause errors for non-configured
trees.

Non-dynamic ports are ensured to be constant by using features already
implemented by the semantic verifier.

The semantic verifier is the last phase of the Jolie interpreter when running
with the \txtl{--check} flag. As we will discuss in the coming section, it is
not practical to perform type-checking of parameters at this point. As a
result, it will not be checked using the \txtl{--check} flag.

\subsection{Additions to the Interpretation Tree Builder and Late Verification}

% Introduction

Interpretation tree building was expanded to support parameters. All other
features required no changes to AST structure, and could simply reuse existing
features.

This section will discuss the technical details of how the Interpretation Tree
Builder (ITB) works. It will cover issues that are specific to this particular
implementation of the Jolie engine.

% What is the ITB and what does it do?

The Interpretation Tree Builder (ITB) is the last step before a Jolie program
is ready to be executed. The ITB works by traversing the AST. The ITB will
collect information from the AST in the form of definitions and inform the AST
about their presence. These definitions include ports, interfaces, types, and
code procedures.

% Processes

Within code procedures the ITB will produce ``processes'' from the AST nodes
which are runnable. It is from processes the meat of code execution is
implemented.  For example, it is from processes everything from addition to
operation calls are implemented.

% When are types ready?

In this phase, the ITB will be gathering complete information about types. Type
building is, for the most part, a one-to-one translation from the AST nodes
into the internal representation. An important job of the ITB with respect to
type building is resolving of type references.  This turns references to a name
into an actual type reference.  Because types aren't built until this phase is
what makes type-checking during semantic verification problematic. The semantic
verifier would essentially have to duplicate the work of the ITB in order to
perform type-checking earlier.  This also means that type-checking cannot be
performed until the ITB is finished.

% State, and how the ITB changes initial state.

The ITB will need to initialize the system with certain values. For example,
the ITB will initialize the \joliel{location} and \joliel{protocol} variables
of output ports. These variables are local only to a single request. At this
point, however, the state of the Jolie program has not yet been initialized.

In Jolie, the state is associated with the current thread of execution. Jolie
uses a separate thread of execution for handling every message. As a result,
every single request gets its own fresh state. All state is copied from the
initialization execution thread. This thread performs work requested by the ITB
followed by the \joliel{init} procedure. Thus right after the ITB is done, and
the initialization execution thread is running, code execution starts. This
includes both user code and loading of any embedded services.

It should be noted that any process which works with state must be run on an
execution thread since it depends on the state being present on the thread.

% Handling of parameter definitions and assignments inside the ITB

The ITB receives two new node types from the configurator and code parser:
parameter assignments, and parameter definitions.

For parameter definitions, the ITB will resolve the types. Ensuring that all
links are properly resolved. This will output a ``processed'' parameter
definition which the ITB will inform the interpreter about.

In a similar fashion, the ITB will process parameter assignments and inform the
interpreter. Type-checking cannot be performed at this point, due to types not
being ready before the ITB finishes. Simply outputting processes for
configuration will make it impossible to perform type-checking before code
execution starts, due to their dependency on state.

The processed assignments contain an ordinarily processed LHS (of the
        assignment), which contains the variable path. The RHS has been
processed into an expression, which can be evaluated into a native Jolie value
on which type-checking can be performed.

% Introduction of new phase and working around lack of state

Because of these constraints a new ``late-checking phase'' has been created.
This phase runs right after the ITB, but before code execution starts. In this
phase, it is not possible to directly touch the init state. It is, however,
    possible to both create values, and evaluate expression. This phase creates
    a synthetic state by creating a new value which will act as the state.  The
    expressions are evaluated against this root and later type-checked. The
    result of evaluating the parameter assignments are then copied into the
    init state, assuming, of course, that type-checking was successful.




\chapter{Jolie Packages}
\section{Introduction}

A Jolie Package is an extension of a Jolie Module. Recall that a Jolie Module
was defined as a directory root, a name, and optionally an entry-point
for the module. A package extends this concept by adding information required
for package management.

A Jolie Package is described by a package manifest. The package manifest is a
JSON file, which is always placed at the root of the package, and must be
called \verb!package.json!. The fixed location allows for the package manager
to easily identify a package. The JSON format was chosen as it plain-text, and
easy to read and write for both humans and machines.

In listing \ref{lst:simple_manifest} we show a simple package manifest. This
manifest showcases the most important features of the manifest. A complete
specification of the package manifest format can be seen in Appendix 1. The
service that this manifest describes is shown in Figure \ref{fig:simple_calc}.

\begin{listing}[H]
\begin{minted}{json}
{
    "name": "calculator",
    "main": "main.ol",
    "description": "A simple calculator service",
    "authors": ["Dan Sebastian Thrane <dathr12@student.sdu.dk>"],
    "license" "MIT",
    "version": "1.0.0",
    "dependencies": [
        { "name": "addition", "version": "1.2.X" },
        { "name": "multiplication", "version": "2.1.0" }
        { "name": "numbers", "version": "1.0.0" }
    ]
}
\end{minted}
\caption{A simple package manifest}
\label{lst:simple_manifest}
\end{listing}

This manifest configures the ongoing example of the calculator system.  Lines
2-3 take care of the module definition, the root is implicitly decided by the
location of the manifest. The remaining attributes, however, are entirely
unique to packages. Some attributes included in the manifest are there for
indexing and discoverability purposes, examples of such attributes are shown in
lines 4-6. The version we will cover shortly, but it is used for package
management. The rest of the manifest describes the dependencies of this
package.



\section{Architecture}
\label{sec:jpm_simple_arch}

% At this point we have covered the responsibilities of JPM. Here we will be
% discussing the actual technical details of JPM. As well as taking a deeper
% look into features provided, and why they are useful.

Before continuing with the packages we'll take a quick de-tour to look at the
overall architecture of JPM. This should help the reader understand in which
context the packages exist. The details of this system will be covered in
Chapter \ref{cha:pkg_mngr}.

The services of the ecosystem around JPM is written in Jolie, using the
features that the Jolie Module System provides, along with the features that
JPM itself provides. Some library services in the ecosystem are written in
Java, using the powerful embedding features of Jolie, which allow certain
foreign languages to be embedded as services.

At the ten-thousand foot view of the architecture, it consists of three
core services, as shown in figure \ref{fig:high_level_arch}:

\begin{enumerate}

\item \textbf{Registry}: Responsible for serving packages known to the
registry.

\item \textbf{JPM}: Provides the back-end of JPM. This includes communication
with one, or more, registries, for example to download packages.

\item \textbf{CLI}: Provides the front-end of JPM. The front-end is responsible
for displaying a user-facing interface, and will communicate with the back-end
to perform the actual work.

\end{enumerate}

\begin{listing}[H]
\begin{minted}{text}
                +-----------+     +-------+     +------------+
                |  jpm-cli  | ==> |  jpm  | ==> |  registry  |
                +-----------+     +-------+     +------------+
\end{minted}
\caption{Ten-thousand foot view of the JPM architecture}
\label{fig:high_level_arch}
\end{listing}


\section{Package Manifest}

In this section we'll cover, in details, the format of Jolie manifests. The
complete specification for the package manifest can be found in Appendix A.
However in this section we'll also cover some of the reasoning behind the
choices found in Appendix A.

\subsection{The {\tt name} Attribute}

The \mintinline{text}{name} attribute uniquely identifies a package. We'll
start this discussion by looking at the constraints imposed on it, from
Appendix A:

\begin{enumerate}
\item
    The name of a package is \emph{not} case-sensitive
\item
    The length of a name is less than 255 characters
\item
    Names may only contain unreserved URI characters (see section 2.3 of
    \href{https://www.ietf.org/rfc/rfc3986.txt}{RFC 3986})
\item
    Names are US-ASCII
\item
    Every registry must only contains a single package with a given name.
\end{enumerate}

Points 1 through 4 all contribute towards a common goal: packages should be
displayable in URLs. This will allow us to more easily create, for example,
a web application where a developer can browse through packages.

Putting some constraints on the package name is also helpful in multiple places
of the package manager. Consider for example a system that might store packages
organized by their name (say for example a caching mechanism). If one package
would be named ``../foo'', then that might accidentally override the contents
of the package named ``foo''. Of course these name constraints are not a
general solution for these types of injection attacks, but it does allow for
several services working with packages to not worry as badly about the package
names.

It should also be noted that the package names are not compatible with the
rules of identifiers. This has affected the design of the Jolie module system.
This is, for example, the reason that strings are used in both configuration
files and module includes, when making references to other modules.

Rule 4 states that package names are unique within a single registry. But this
constraint actually goes even further, since these packages directly extend
from the modules, we cannot have a system containing two packages named the
same thing, even if they reside in two different registries. Generally we make
no assumptions that registries communicate with each other, for this reason it
is perfectly allowed that two distinct registries may contain name clashes, but
any single package must not have a dependency which clashes either with another
dependency, or the package who has the dependency itself.

\subsection{The Meta Data Attributes ({\tt license, authors, description})}

These attributes are mostly used for search results. They allow for a developer
seeking information about a package, to quickly gather the most important
information.

License identifiers are validated against a pre-existing list of license
identifiers. This is done to ensure that typos are not made in license names.
The identifiers used are from the SPDX license list, which contains a list of
commonly found licenses. Currently only the license identifiers are supported,
this could be extended to use SPDX License Expression Syntax, which allows for
references to custom licenses.

\subsection{The {\tt version} Attribute}

% source: Semantic Versioning 2.0.0 - Tom Preston-Werner - http://semver.org/

The \mintinline{text}{version} attribute used by packages declare which version
of the source code the manifest describes. This field uses the Semantic
Versioning (SemVer) specification for version numbers. It is a requirement to
use SemVer versions, if a version is specified, which is not a valid SemVer
version, then validation of the manifest will fail. Enforcing this standard
allows for several other features, which will be covered in a moment.

SemVer version numbers consists of three fields, all integers, those being the
\emph{major}, \emph{minor}, and \emph{patch} version numbers. Additionally a
version may have a pre-release label or additional meta data. Labels and build
meta-data won't be covered, as they are not that important for our use-case.
The (simplified) syntax of SemVer version number is shown in Listing
\ref{lst:semver_syntax}.

\begin{listing}[H]
\begin{minted}{text}
<major>.<minor>.<patch>[-<label>][+<meta-data>]
\end{minted}
\caption{Simplified syntax of a SemVer version number}
\label{lst:semver_syntax}
\end{listing}

Precedence of version numbers are determined by comparing each unit numerically
in the order of major, minor, and finally patch. The first differentiating
units determine the precedence. Thus, for example: 1.0.0 $<$ 1.0.1 $<$ 1.1.0
$<$ 1.10.0 $<$ 2.0.0.

Semantic Versioning dictates rules for how version numbers should change,
depending on how the public API changes. The rules can be summarized as:

\begin{enumerate}

    \item The patch version number \emph{must} be incremented, if only
        backwards compatible bug fixes are introduced.

    \item The minor version number \emph{must} be incremented, if backwards
        compatible functionality is introduced.

    \item The major version number \emph{must} be incremented, if incompatible
        changes are introduced.

\end{enumerate}

Having a major version of 0, indicates that the product is not yet stable. At
this point normal rules for incrementing version numbers do not apply. As a
result version 1.0.0 indicates the first public API.

Semantic Versioning has become a commonly used standard across many different
package managers, examples include
Cargo\footnote{\url{http://doc.crates.io/manifest.html}} and
NPM\footnote{\url{https://docs.npmjs.com/getting-started/using-a-package.json}}.

TODO Highlight benefits of using semantic versioning.

% source: https://docs.npmjs.com/misc/semver

For dependencies the version field allows for a bit more flexibility using
SemVer expressions, also known as SemVer ranges. These expressions allows the
developer more freedom in which version to use. In short the expressions allows
the developer to both whitelist certain ranges for use, and blacklist certain
ranges for use.

This provides both a convenience factor, i.e. the developer can express that
she doesn't care which specific version is used, as long as it has feature X
(released in some known version).

This feature also provides benefits for when the dependency tree is to be
calculated. Consider, for example, two dependencies A and B which both depend
on C with the same level of functionality, which was released in version 1.2.0.
Assuming that expressions were not allowed, these two services would have to
describe their dependency in an exact version. Service A might have been
written at a time when 1.2.3 of service C was the newest, while B was written a
bit later and is thus depending on version 1.2.6. The problem is now that the
client depending on A, and B, cannot do so, because A and B have conflicting
requirements, despite the fact that they both require the same level
functionality. With SemVer expressions these services may more accurately
describe what they actually depend on, service A might state that it depends on
">=1.2.0", while B is a bit more conservative and wants exactly version 1.2.6.
In this case the package manager will be able to choose version 1.2.6, since
this version fulfills the requirements of all dependencies.

Allowing for SemVer expressions to be used, is however not without its
drawbacks. In Section \ref{sec:lockfiles} we cover one of these problems, and
its solution.

\subsection{The TODO Name Attributes ({\tt events}, {\tt main})}

These attributes control aspects of the package related to the execution of
scripts, by the package manager, related to the package.

The \mintinline{text}{main} attribute controls the entry-point of the module.
This is used both in informing the Jolie Engine where the entry-point is
located (Section TODO). The attribute is also used for starting the package, we
discuss this in Section TODO.

The \mintinline{text}{events} attribute defines a series of lifetime hooks,
which are scripts that are run when specific lifetime events occur, this is
described in Section TODO.

\subsection{The {\tt registries} Attribute}

The package manager may communicate with one or more registries, when
performing its work. This may either be to install dependencies, or it could be
to perform some of the more organizational commands, such as creating an
account.

For this the package manager will need to know about the different registries.
This is what the \mintinline{text}{registries} attribute describes, it is an
array containing definitions for all known registries (known to the package).

A registry is defined by a unique (within the package) name, and a location (a
URI following the same rules that Jolie does for its port locations). Listing
\ref{lst:registry_definition} shows an example.

\begin{listing}[H]

\begin{minted}{json}
{
    "name": "my-registry",
    "location": "socket://registry.example.com:9999"
}
\end{minted}

\caption{A registry named \mintinline{text}{my-registry} being hosted at
    \mintinline{text}{registry.example.com} running on port 9999}

\label{lst:registry_definition}

\end{listing}

Every package manifest has an implicit entry which is named ``public''. This
entry points to the public registry. Every command which needs to use a
registry, will optionally receive a registry name. If this name is specified,
it will look within the \mintinline{text}{registries} attribute. If no name is
specified, the ``public'' entry will be used.

TODO Should we talk about why it is good to have multiple registries here?

\section{Dependencies}

A JPM dependency describes a dependency on another JPM package. Since JPM
packages are extensions of Jolie modules, what we're really describing is a
code dependency.

There is an important, but subtle, distinction between the dependencies of a
Jolie service and the dependencies listed in a JPM package manifest. For the
sake of clarity, we will refer to the first one as service dependencies, and
the latter as package dependencies.

Consider, the Jolie service shown in Listing \ref{lst:jolie_dep_service}.  For
this service we can define one explicit service dependency on the
\txtl{Example} service and one implicit service dependency on the client
calling the \txtl{Self} service.

However given this service, encapsulated in a package, it would have no package
dependencies. For JPM a package dependency is a dependency on the code of
another module. In this example, however, we have no dependency on external
code.

\begin{listing}[H]
\begin{minted}{jolie}
// service.ol
interface ExampleInterface { RequestResponse: example(undefined)(undefined) }

outputPort Example {
    // ...
    Interfaces: ExampleInterface
}

interface SelfInterface { RequestResponse: hello(undefined)(undefined) }

inputPort Self {
    // ...
    Interface: SelfInterface
}

main {
    [hello(request)(response) {
        // ...
        example@Example(exampleRequest)(exampleResponse);
        // ...
    }]
}
\end{minted}
\caption{A simple Jolie service served defined in a single file
    (\txtl{service.ol})}
\label{lst:jolie_dep_service}
\end{listing}

If we change the previous example, such that the interface of \txtl{Example} is
defined in another package named \joliel{"example"}, then we could define a
package dependency on the \joliel{"example"} package. This is shown in Listing
\ref{lst:jolie_dep_package}. This service dependencies in this example remain
unchanged.

\begin{listing}[H]
\begin{minted}{javascript}
/* /example/package.json */
{
    "name": "example",
    "version": "1.0.0"
    /* remaining fields omitted */
}

/* /self/package.json */
{
    "name": "self",
    "dependencies": [{ "name": "example", "version": "1.0.0" }]
    /* remaining fields omitted */
}
\end{minted}

\begin{minted}{jolie}
/* /example/interface.iol */
interface ExampleInterface { RequestResponse: example(undefined)(undefined) }

/* /self/main.ol */
include "interface.iol" from "example"

outputPort Example {
    // ...
    Interfaces: ExampleInterface
}

// input port and interface of Self omitted

main {
    [hello(request)(response) {
        // ...
        example@Example(exampleRequest)(exampleResponse);
        // ...
    }]
}
\end{minted}
\caption{A package dependency is a dependency on the code of another package}
\label{lst:jolie_dep_package}
\end{listing}

A package dependency can be described by the 3-tuple $(name, version,
registry)$, as covered the manifest will implicitly set the registry to
\joliel{"public"} if not specified.

The dependency tree of a JPM package is build starting at the package itself.
The algorithm for calculating the dependency tree is shown, in pseudo-jolie, in
Listing \ref{lst:dep_tree}. Some details, like lockfiles and data marshalling,
are left out. Error handling is illustrated by the \txtl{Error}
procedure. In reality the error handling is slightly more advanced, but
not relevant to our discussion.

\begin{coded}
\begin{minted}{jolie}
define DependencyTree {
    // "manifest" is set at entry to this procedure
    dependencyStack << manifest.dependencies;
    currDependency -> dependencyStack[0];

    while (#dependencyStack > 0) {
        RegistrySetLocation; // Registry now points to origin of dependency

        // Lookup version information from the registry
        getPackageInfo@Registry(currDependency.name)(packageListing);
        if (#packageListing.results == 0) { Error };

        resolved -> resolvedDependencies.(currDependency.name);
        if (is_defined(resolved)) {
            // Check if 'resolved' matches requirement of 'currDependency'
            satisfies@SemVer({
                .version = resolved.version, .range = currDependency.version
            })(versionSatisfied);
            if (!versionSatisfied) { Error };
        } else {
            allVersions -> /* extract versions from packageListing.results */;

            // Find the best match for our version expression
            sortRequest.versions -> allVersions;
            sortRequest.satisfying = currDependency.version;
            sort@SemVer(sortRequest)(sortedVersions);
            if (#sortedVersions.versions == 0) { Error };

            // Insert resolved dependency
            with (information) {
                .version << sortedVersions.versions
                    [#sortedVersions.versions - 1];
                .registryLocation = Registry.location;
                .registryName = registryName
            };
            resolved << information;

            // Insert dependencies of this dependency on the stack
            dependenciesRequest.packageName = name;
            dependenciesRequest.version << resolved.version;
            getDependencies@Registry(dependenciesRequest)(depsResponse);
            for (i = 0, i < #depsResponse.dependencies, i++) {
                item << depsResponse.dependencies[i];
                dependencyStack[#dependencyStack] << item
            }
        }
    }
}
\end{minted}
\caption{Pseudo-Jolie code for calculating a dependency tree}
\label{lst:dep_tree}
\end{coded}

The algorithm keeps two primary data-structures: the \joliel{dependencyStack}
and the \joliel{resolved} versions. The algorithm starts by inserting all the
listed dependencies from the \joliel{manifest} onto the
\joliel{dependencyStack}, this is done in lines 2-5. The algorithm then enters
the main loop (line 7).

This loop will pop a dependency off the stack and start the process of
resolving the dependency to a known version. Recall that dependencies use
SemVer expressions in their version field. We need to resolve this field to a
particular version. In this process we will also confirm that the registry has
knowledge of this version.

This process starts by asking for a complete version listing for a package
matching the name of our dependency (lines 8-12). At this point we can be
certain that the registry at least know of a package by that name. We then
check if a previous dependency already has resolved this dependency.
The \joliel{resolved} dependencies data structure is a dictionary, mapping
a package name to a resolved version. There are two important takeaways from
this decision:

\begin{enumerate}

\item \textbf{A package cannot have two (or more) identically named
    dependencies.} Even though it is technically possible to have two
    identically named dependencies, as long as they do not originate from the
    same registry, this is still disallowed. The reason for this lies in the
    module implementation, where modules are uniquely identified only by their
    name. As a result we would have to either identify all packages by their
    name and registry, or simply disallow two identically named dependencies.
    We chose the latter to simplify the use of packages in code.

\item \textbf{A package cannot depend on two (or more) different versions of a
    package}. This is similar to the first point, but made worse by the lack
    of any namespacing in Jolie. Because of this, if two packages were to
    provide any two constructs which clash in name, the Jolie engine would be
    unable to distinguish them.

\end{enumerate}

If a dependency has not been resolved (lines 20-46) we will start the resolving
phase. The package manager will \emph{always} select the newest version
available in the registry which matches the SemVer expression as defined by the
dependency (lines 21-27). When the package has been resolved its dependencies
are looked up (by contacting the registry, we don't have the actual manifest).
These dependencies are then added onto the stack (lines 38-46).

If the dependency has already been resolved, we must check if this dependency
works under these conditions, by matching its SemVer expression to the resolved
dependency. If this is possible we can continue. \emph{The package manager
    makes no additional attempts to resolve the constraints imposed by the
        manifests}.

These constraints mean that there plenty of dependency trees can be
constructed, such that the greedy approach of selecting the newest version,
    will lead to conflicts. This approach was chosen for its simplicity. More
    advanced approaches were not explored. The greedy approach will most likely
    also be significantly more efficient. At least one other package manager
    was also found to be using this exact same approach to constraint
    solving\cite{YARNA}.

\section{Lock Files}
\label{sec:lockfiles}

% https://www.martinfowler.com/articles/continuousIntegration.html
% https://en.wikipedia.org/wiki/Continuous_integration#cite_note-2

Lock files are a feature designed to solve some of the drawbacks associated
with allowing SemVer expressions in the version field of dependencies. To see
how SemVer expressions pose a problem, we must first inspect the typical
engineering process in which it is used. This feature is by no means original,
prior work includes: Cargo\autocite{CRAA}, and Yarn\autocite{YARNB}.

Continuous integration (CI) is a software development process\autocite{MFCI}.
This process is used to avoid the so called ``integration hell'', in which the
time taken to integrate, or merge, the changes from multiple developers into a
single track, takes longer than developing the individual features.  The
process works by, multiple times a day, integrating all these changes, this way
failure in integration is found much earlier. We'll quickly summarize one
possible implementation of CI here. This is not the only way to do it, but
summarizes the most important parts, which we'll use to highlight potential
problems with SemVer expressions.

For this process to work efficiently it frequently uses automated tests and
builds, in order to perform this integration. The automated tests and builds
are then performed by a dedicated server. The server will most likely pull the
source code from a single repository, which represents the most up-to-date
track of development.

When new features are developed, the developer will pull the most up-to-date
source code, and develop the feature locally with this. This includes
performing all of the builds and tests. Thus when the feature is pushed back
onto the source control all of these tests should still pass. The role of the
CI server is then to verify that the process is upheld.

How the build and testing is performed obviously depend on which technologies
are used. A typical JPM based project might perform the following steps:

\begin{enumerate}
\item Pull source code from version control
\item Perform the build
    \begin{enumerate}
        \item Download JPM dependencies (i.e. \verb!jpm install!)
        \item Potentially build other necessary artifacts (such as Java
                libraries)
    \end{enumerate}
\item Run tests
\end{enumerate}

It is in step 2(a) the need for lock files arises, the problem is most easily
demonstrated through an example.

Consider the following scenario: a developer wishes to develop a new feature
which uses package X of at least version 2.1.0. For this reason the developer
adds the following line to the dependencies section: \mintinline{json}{{
"name": "X", "version": ">=2.1.0" }}. Afterwards the developer performs an
install of this package, at the moment this dependency to version 2.1.3. At
this point the developer continuous with developing the feature, without
having to perform any additional installs of the package. Once the feature
is developed, the developer performs all of the tests, and see that they
pass.

At this point the code is pushed onto the source control system, where it is
picked up by the CI server for testing. The CI server will then perform all of
the steps listed above, including installing the dependencies anew. However,
since the initial install by the developer had been performed a new version
has been released, say version 2.3.0. This version still fulfills the
expression listed in the dependency, and is hence chosen, since this is the
best fitting package for this dependency. Remember, the CI server has no way
of knowing which package was originally installed, it only has the package
manifest to work with.

Installing a different version of the package may however lead to the tests no
longer working! As a result, we cannot guarantee what works on a developer's
machine, necessarily works on any other machines that might try to run it.
As a result, some might choose to never use SemVer expressions to avoid this
pitfall.

Lock files offer a compromise between these two extremes. A lock file contains
the exact versions that every dependency was resolved to. The lock file should
be put into source control, thus when another machine attempts to perform the
build the package manager will know which exact version should be used for the
build.

The convenience factor of SemVer expressions is present. When we perform the
initial install, we might not care so which version to use, and the package
manager will still be able to just pick the best version for us.  Additionally
the package manager can provide upgrade scripts, which uses these expressions
as guidelines.

The lock files is placed in the package directory, and is called
\verb!jpm_lock.json!. Listing \ref{lst:lock_file} shows an example of a lock
file, which would have been generated for the scenario above.

The lock files of a dependency is \emph{not} used by a client package. As a
result library developers should still be careful not specifying too wide
expressions. This feature is only intended to guarantee that a package will
have the same dependencies regardless of which machine it runs on. It is not
intended as a replacement for specifying exact versions. If a package requires
a specific version, then that should be reflected in the package manifest.

\begin{listing}[H]
\begin{minted}{json}
{
    "_note": "Auto-generated file notice",
    "locked": {
        "X@>=2.1.0/RegistryName": {
            "resolved": "2.1.3",
            "checksum": "..."
        }
    }
}
\end{minted}

\caption{A lock file showing that the dependency for package X of at least
    version 2.1.0 has been resolved to version 2.1.3}

\label{lst:lock_file}
\end{listing}

\section{Lifetime Hooks}
\label{sec:lifetime_hooks}

A lifetime hook is a script that ``hooks'' onto certain lifetime events that
occur during normal use of the package manager. These scripts allows the
package developer to augment the package manager, and potentially affect the
work it performs.

This is a feature which is available in many other types of software, for
example the popular version control system Git provides such a
feature\autocite{chacon2009pro}.

All hooks are performed on the client-side, and never on the server side (i.e.
on registries). Only the hooks of the current package will be run, the
hooks of a dependency will not be executed. Hooks that run before a certain
action, denoted by a \verb!pre-! prefix, may stop the corresponding command
from being executed.  These scripts may do so by returning a non-zero exit
code.

These features become useful for augmenting the workflow with in-house
conventions. This may for example include ensuring that, for example, tests
pass before publishing a new version. Using a simple process interface allows
for the user to write these scripts in whichever technology they choose.

The hooks are defined in the ``events'' section of the package manifest.
Listing \ref{lst:hook} shows an example hook, which runs before publishing
which ensure that all tests pass.

\begin{listing}[H]
\begin{minted}{javascript}
{
    // ...
    "events": {
        "pre-publish": "./run_tests.sh"
    }
}
\end{minted}

\caption{Defining a lifetime hook, which runs the script
    \mintinline{text}{run_tests.sh} before publishing the package to the
        registry}

\label{lst:hook}

\end{listing}

The following are the hooks that JPM supports:

\begin{itemize}

\item \verb!pre-start!: Script is run right before the \verb!start! command is
    invoked on a service.

\item \verb!post-start!: Run right before the \verb!start! command
    terminates. There is no guarantee that this script is run (i.e. if the
    service was forcefully terminated).

\item \verb!pre-install!: Runs before the \verb!install! command is invoked.

\item \verb!post-install!: Runs after the \verb!install! command has finished.

\item \verb!pre-publish!: Runs before the \verb!publish! command is invoked.

\item \verb!post-publish!: Runs after the \verb!publish! command has finished.

\end{itemize}

\section{The {\tt .pkg} Format} \label{sec:pkg}

The \verb!.pkg! file format, is the format that JPM uses for distribution of
packages. Simply the file format contains a Jolie package zipped up into a
single file.

Currently the contents of the zip file maps directly with the underlying
package. The format is however open to extension, by allowing for new files to
be introduced that can be used in various extensions. One such extension could
for example be native code signing, which we discuss in more details in Section
TODO. An obvious implementation strategy would involve adding the necesary
files directly into the ZIP archive.

Extending the ZIP file format is a quite common occourence. Espescially when
shipping what is essentially a collection of files. Popular examples of other
file formats using this approach includes the JAR\footnote{Package file format
    typically used for aggregating many Java class files together} format, and
    Office Open XML\footnote{Developed by Microsoft, used in their office
        applications}.



\chapter{Package Manager}
\label{cha:pkg_mngr}
\section{Introduction}

The Jolie Package Manager is build to facilitate the use of packages, which
were introduced in the previous chapter. The package manager provides a variety
of tools, these can roughly be divided into three categories:

The command line application serves as the user interface to JPM. The
application is, perhaps not unsurprisingly, named \mintinline{text}{jpm}. The
tool will be used in several examples.

\begin{enumerate}
\item \textbf{Package management}
    \begin{itemize}
        \item Installing packages (Section \ref{sec:cache})
        \item Publishing packages (Section TODO Registry-database)
        \item Upgrading packages (Section \ref{sec:versions} and
                \ref{sec:lockfiles})
        \item Searching for packages (Section TODO Registry-database)
    \end{itemize}
\item \textbf{Account management (Section \ref{sec:security})}
    \begin{itemize}
        \item Login/logout with registries
        \item User management
        \item Team management
    \end{itemize}
\item \textbf{Helper scripts}
    \begin{itemize}
        \item Creating a new package (Section TODO New helper scripts)
        \item Starting a package (Section TODO New helper scripts)
        \item Interacting with the cache (Section \ref{sec:cache})
    \end{itemize}
\end{enumerate}

These responsibilities are services by the package manager ecosystem which were
briefly covered in Section \ref{sec:jpm_simple_arch}. Many of the features this
ecosystem provides were also covered in Chapter \ref{cha:packages}. In this
chapter we will cover the final details of this ecosystem.


\section{The Command Line Interface}

The command line application is responsible for displaying a more user friendly
interface to inner workings of JPM. The tool will perform almost not work by
itself, but will instead delegate this to the back end (See Section TODO).

When first running the tool, the user will be welcomed with the following
message:

\begin{minted}{text}
JPM - The Jolie Package Manager
Version 1.0.0

Usage: jpm <COMMAND> <COMMAND-ARGUMENTS>

Command specific help: jpm help <COMMAND>

Available commands:
-------------------
init           Initializes a repository
search         Searches repositories for a package
install        Install dependencies
publish        Publish this package
start          Start this package.

[ Remaining commands removed from snippet ]
\end{minted}

As clearly visible form this snippet, for the tool to do any work we must first
give it something to do via a command. The commands that JPM understands almost
directly mirror the functionality provided by the back end. To use JPM to
create a new package, the user must simply use the \mintinline{text}{init}
command. This will display a prompt, guiding the user through the mandatory
field, and automatically create a package with the required structure. This
is shown in Listing \ref{lst:jpm_init}.

\begin{listing}[H]
\begin{minted}{text}
$ jpm init
Package name
------------
> my-package

Package description
-------------------
> This is my package

Author: [Format: name <email> (homepage)]
-----------------------------------------
> Dan Sebastian Thrane <dathr12@student.sdu.dk> (github.com/DanThrane)

Private package? [Y/n]
----------------------
> n

$ cat my-package/package.json | json name
my-package
\end{minted}

\caption{The \mintinline{text}{jpm} tool provides a user interface for common
    tasks. In this example, creating a new package.}

\label{lst:jpm_init}

\end{listing}

\subsection{Internal Organization and Deployment}

The command line tool delegates most of the work to the back end service (here
called jpm). Figure \ref{fig:cli_arch} shows the architecture from the CLI's
point of view. Most noteably is the callback server, this server is responsible
for receiving information about events that occur in the back end. These are
primarily used to communicate progress, this is especially useful for long
running processes, such as downloading dependencies. The back end will in these
cases these events to a callback server, which can then choose to display
information about this event. The need for a separate service comes mostly from
a limitation in Jolie. In Jolie all communication follows either a one-way, or
request-response communication pattern. As a result of this, it isn't possible
for JPM to send back information while a request is being processed. To handle
this the front end (in this case \mintinline{text}{jpm-cli}) will inform the
back end of where it should send events.

\begin{listing}[H]
\begin{minted}{text}
                 +-------------------------------------+
                 | terminal (user)                     |
                 +-------------------------------------+
                                   |
                                   v
                 +-------------------------------------+
+------------+   |            embeds                   |
| console-ui |<--| jpm-cli ----------------+           |
+------------+   |                         |           |
+------------+   |                         v           |
| arg-parser |<--|                   +----------------+|
+------------+   |                   | callback       ||
                 |                   +----------------+|
                 +--------------------------^----------+
                                   |        |
                                   |        |
                                   v        |
                 +-------------------------------------+
                 |                jpm                  |
                 +-------------------------------------+
                     |         |         |         |
                     v         v         v         v
                 +-------+ +-------+ +-------+ +-------+
                 | reg-a | | reg-b | | reg-c | | reg-d |
                 +-------+ +-------+ +-------+ +-------+
\end{minted}
\caption{The system architecture from the CLI's point of view}
\label{fig:cli_arch}
\end{listing}

Quite a lot of services exists on the client side, notable
\mintinline{text}{jpm-cli} and \mintinline{text}{jpm}. From a usability
perspective this is less optimal. For this reason the \mintinline{text}{jpm}
binary, which ships with the package manager, will use a deployment which
embeds all of the core services together. Future versions of the package
manager, may wish to optionally spin up the required services, and run them as
daemons\footnote{A daemon is a background process running on a computer}. This
will remove quite a bit of the overhead associated with spinning up the JVM
(required for the Jolie engine) and every embedded service. Such a practice is
done by similar tools, as an example the JVM based build tool Gradle provides
such a daemon\autocite{GRAA}.  This conversion would be relatively straight
forward, since the server architecture required by the daemon is already
implemented.

\section{The JPM Back End}

The back end service is responsible for performing most of the heavy lifting,
for the commands that are send to it from the front end. Some of the jobs
are performed directly by this service, while quite a few others are
delegated.

The most important partner the back end service has is the registry service.
The back end will communicate with one or more registry services.

The JPM back end collaborates with a few external services, these are
illustrated in Table \ref{tab:collab_jpm}. The work performed by the back-end
has mostly already been discussed or will be discussed (see Table).

\begin{table}[H]
\centering


\begin{longtable}[c]{@{}lll@{}}
\toprule
\begin{minipage}[b]{0.17\columnwidth}\raggedright\strut
Service Name
\strut\end{minipage} &
\begin{minipage}[b]{0.60\columnwidth}\raggedright\strut
Responsibilities
\strut\end{minipage} &
\begin{minipage}[b]{0.14\columnwidth}\raggedright\strut
Sections
\strut\end{minipage}\tabularnewline
\midrule
\endhead
\begin{minipage}[t]{0.17\columnwidth}\raggedright\strut
\txtl{registry}
\strut\end{minipage} &
\begin{minipage}[t]{0.60\columnwidth}\raggedright\strut
Responsible for storing packages
\strut\end{minipage} &
\begin{minipage}[t]{0.14\columnwidth}\raggedright\strut
\ref{sec:registry}
\strut\end{minipage}\tabularnewline
\begin{minipage}[t]{0.17\columnwidth}\raggedright\strut
\txtl{execution}
\strut\end{minipage} &
\begin{minipage}[t]{0.60\columnwidth}\raggedright\strut
Used for starting packages and executing life-time hooks
\strut\end{minipage} &
\begin{minipage}[t]{0.14\columnwidth}\raggedright\strut
N/A
\strut\end{minipage}\tabularnewline
\begin{minipage}[t]{0.17\columnwidth}\raggedright\strut
\txtl{lockfiles}
\strut\end{minipage} &
\begin{minipage}[t]{0.60\columnwidth}\raggedright\strut
Handles the lockfiles feature
\strut\end{minipage} &
\begin{minipage}[t]{0.14\columnwidth}\raggedright\strut
\ref{sec:lockfiles}
\strut\end{minipage}\tabularnewline
\begin{minipage}[t]{0.17\columnwidth}\raggedright\strut
\txtl{semver}
\strut\end{minipage} &
\begin{minipage}[t]{0.60\columnwidth}\raggedright\strut
Library package for semantic versioning
\strut\end{minipage} &
\begin{minipage}[t]{0.14\columnwidth}\raggedright\strut
\ref{sec:versions}
\strut\end{minipage}\tabularnewline
\begin{minipage}[t]{0.17\columnwidth}\raggedright\strut
\txtl{jpm-downloader}
\strut\end{minipage} &
\begin{minipage}[t]{0.60\columnwidth}\raggedright\strut
Acts as a cached downloader and package installer
\strut\end{minipage} &
\begin{minipage}[t]{0.14\columnwidth}\raggedright\strut
\ref{sec:cache}
\strut\end{minipage}\tabularnewline
\begin{minipage}[t]{0.17\columnwidth}\raggedright\strut
\txtl{packages}
\strut\end{minipage} &
\begin{minipage}[t]{0.60\columnwidth}\raggedright\strut
Validator for the package format
\strut\end{minipage} &
\begin{minipage}[t]{0.14\columnwidth}\raggedright\strut
\ref{sec:manifest}
\strut\end{minipage}\tabularnewline
\begin{minipage}[t]{0.17\columnwidth}\raggedright\strut
\txtl{pkg}
\strut\end{minipage} &
\begin{minipage}[t]{0.60\columnwidth}\raggedright\strut
Creates the binary format used for distribution
\strut\end{minipage} &
\begin{minipage}[t]{0.14\columnwidth}\raggedright\strut
\ref{sec:pkg}
\strut\end{minipage}\tabularnewline
\begin{minipage}[t]{0.17\columnwidth}\raggedright\strut
\txtl{callback}
\strut\end{minipage} &
\begin{minipage}[t]{0.60\columnwidth}\raggedright\strut
Callback service. JPM calls this about events
\strut\end{minipage} &
\begin{minipage}[t]{0.14\columnwidth}\raggedright\strut
\ref{sec:cli}
\strut\end{minipage}\tabularnewline
\bottomrule
\end{longtable}

\caption{Collaborators with the JPM backend}
\label{tab:collab_jpm}
\end{table}


\section{Registry}
\label{sec:registry}

In this section we will briefly cover the responsibilities of JPM registries.
For the most part we will cover the specific features it provides in other
section and make references to them, when relevant.

The primary responsibility of a JPM registry is to manage packages and serve
them to clients who request them. When a developer publishes a packages then
this is done to a particular package registry. When a developer downloads a
package this is done from a registry. Almost anytime the \txtl{jpm} tool needs
to do anything it must contact a registry for information.

Many registries can exist in the JPM ecosystem. A default registry, the public
registry\footnote{As of writing this thesis, the module system and package
    system has not yet been merged into the Jolie language. As a result, when
    referring to the public registry, we're referring to a registry which might
    exist in the future. The software required to run a registry is, however,
    implemented, as described in this thesis. A public registry does
    not exist yet.}, is made for sharing open-source packages with the Jolie
community. Operations in the \txtl{jpm} tool that needs to speak to a registry
will default to this registry.

Having multiple registries is beneficial for example if an organization wishes
to share packages, but only wishes to share them internally within the
organization. The JPM registry, and all other packages in the JPM ecosystem,
are JPM packages. They are available from the public registry. This makes it
relatively easy to host your own registry, since the process of this is the
exact same as any other Jolie package.

The \registry package also provides user and team management. Different rights
may be assigned to these users, which can control who may perform certain
actions. This is needed for controlling who's allowed to publish updates for a
package. This also allows for a registry to be configured to only allow users
from within an organization to use it. These features are covered by the
\security package which is described in Section \ref{sec:security}.

Keeping track of package information is delegated to the \regdb package. We
cover this package in Section \ref{sec:regdb}.

\section{The Package Database ( \regdb)}
\label{sec:regdb}

The \regdb package is responsible for storing information about packages. In
this section we will cover in depth how this service works and how it is used.

This package provides storage, and querying facilities for package information.
This package does not provide any storage of actual packages. It is left to the
client of this package to perform this. This allows freedom in how contents is
stored, which may change between clients.

This package is consumed by both the \registry and the \cache packages. How
these packages communicate with the \regdb package will be covered in their
respective packages. The \regdb package essentially acts as a front-end to a
SQL database. We won't cover the technical details of how this works, but for
the remainder of this section we will cover what is stored and which facilities
this package provides.

The core data stored in the database is shown in Figure
\ref{fig:registry_database}. The core database model consists of three tables:
\txtl{package}, \txtl{package_dependency}, \txtl{package_versions}. The
\txtl{package} table describes a particular known package. Using any particular
entry in the \txtl{package} table, we may find all known versions of this
package. From a particular version of the package we may find all its known
dependencies.

\begin{figure}[H]
    \centering
    \includegraphics[width=1.0\textwidth]{pictures/regdb.png}
    \caption{The core database model used in the \regdb package}
    \label{fig:registry_database}
\end{figure}

The \txtl{origin} fields which appear in all three tables reference the origin
registry. This is used for handling both cross-registry dependencies, but also
for storing information about packages from different registries. The latter is
useful for cache-like services, which collect information from many different
registries.

Any particular package version also contains a \txtl{checksum} field. This
field can contain any generic checksum which can be serialized as text, that is
the actual checksum algorithm is defined by the client.

The remainder of the fields stored in the database simply mirror the most
important fields of the package manifests. Certain fields from the package
manifests, like the lifetime hooks, are left out of this database model. The
reason for this being lack of interest by clients. The remaining fields, like
name, description, license, are all of interest when it comes to
discoverability of packages.

Table \ref{tab:regdb_operations} shows the operations provided by the \regdb
package. Generally the operations provided by this package for querying
operations operate on a key matching either the \txtl{package} or
\txtl{package_versions} table. Operations that update the tables usually
operate on package manifests (typically provided by the \txtl{packages}
package).

\begin{table}[H]

\begin{longtable}[c]{@{}ll@{}}
\toprule
\begin{minipage}[b]{0.35\columnwidth}\raggedright\strut
Operation
\strut\end{minipage} &
\begin{minipage}[b]{0.59\columnwidth}\raggedright\strut
Description
\strut\end{minipage}\tabularnewline
\midrule
\endhead
\begin{minipage}[t]{0.35\columnwidth}\raggedright\strut
\txtl{query}
\strut\end{minipage} &
\begin{minipage}[t]{0.59\columnwidth}\raggedright\strut
Searches for packages matching a text query
\strut\end{minipage}\tabularnewline
\begin{minipage}[t]{0.35\columnwidth}\raggedright\strut
\txtl{checkIfPackageExists}
\strut\end{minipage} &
\begin{minipage}[t]{0.59\columnwidth}\raggedright\strut
Check if any package exists with name origin
\strut\end{minipage}\tabularnewline
\begin{minipage}[t]{0.35\columnwidth}\raggedright\strut
\txtl{getInfoAboutPackage}
\strut\end{minipage} &
\begin{minipage}[t]{0.59\columnwidth}\raggedright\strut
Lists all versions of a package matching name, origin, and possibly
version
\strut\end{minipage}\tabularnewline
\begin{minipage}[t]{0.35\columnwidth}\raggedright\strut
\txtl{compareWithNewest}
\strut\end{minipage} &
\begin{minipage}[t]{0.59\columnwidth}\raggedright\strut
Checks if a package manifest is newer than what is known by this
registry
\strut\end{minipage}\tabularnewline
\begin{minipage}[t]{0.35\columnwidth}\raggedright\strut
\txtl{insertNewPackage}
\strut\end{minipage} &
\begin{minipage}[t]{0.59\columnwidth}\raggedright\strut
Inserts a new version of a package manifest
\strut\end{minipage}\tabularnewline
\begin{minipage}[t]{0.35\columnwidth}\raggedright\strut
\txtl{createPackage}
\strut\end{minipage} &
\begin{minipage}[t]{0.59\columnwidth}\raggedright\strut
Creates a package with name and origin
\strut\end{minipage}\tabularnewline
\begin{minipage}[t]{0.35\columnwidth}\raggedright\strut
\txtl{getDependencies}
\strut\end{minipage} &
\begin{minipage}[t]{0.59\columnwidth}\raggedright\strut
Returns a list of dependencies for a package (name, origin, version)
\strut\end{minipage}\tabularnewline
\bottomrule
\end{longtable}

\caption{Operations provided by the \regdb package. Operations named changed
slightly for formatting reasons}

\label{tab:regdb_operations}

\end{table}

\section{The JPM Cache}
\label{sec:cache}

The \cache of JPM is responsible for downloading packages, and maintaining the
packages in a local cache.

The \cache acts as a database of known packages. This database contains
both the information needed to know about the database, as well as the binary
containing the package data.

The \cache consists of previously downloaded packages, along with a
locally kept database which matches the type that the \registry uses. The
database kept also track the origin \registry, since the \cache will
potentially download from multiple. The key used for this is the location URI
of each \registry. The name cannot be used, since these differ between each
package.

The primary operation that the \cache exposes is the \txtl{installDependency}
operation. This operation will install a dependency directly into a package.
JPM assumes that dependencies are installed under
\txtl{<pkgRoot>/jpm_packages/<depName>}. Having packages stored under a common
folder makes it easy to locate. This is, for example, useful when needing to
filter out the packages. Common places for this would be version control, or
from JPM when publishing a package.

The \cache will first attempt to retrieve the package from its local database,
if not found it will forward the request to the registry. In order to
forward the request, the operation is fed both the registry and
authentication token.

When a new package is received from a \registry, two parts will be returned:
the binary itself, along with the checksum that the \registry believes the
package to have. Checksums are both checked when initially receiving the
package, but also at every install from the cache. This helps prevent initial
corruption and further protects from the cache itself being corrupted.

The \cache will maintain a \regdb just like the \registry component does. When
new entries are loaded into the \cache, it will updated the \regdb in a similar
fashion. As a result the \cache is technically capable of performing the exact
same responsibilities as the \registry. This could for example allow for fully
offline downloads assuming the data is already loaded into the cache.


% Status: Draft

\section{Authorization}

In this section we need to discuss:

\begin{itemize}
\item Access control matrix, and how that stuff works.
\item User registration and best practices for handling passwords
\end{itemize}


\section{Authentication Tokens}

In Section \ref{sec:authorization} we covered the authorization framework that
JPM uses. When a user successfully authenticates the authorization framework
will provide an authentication token proving the user's identity.

For user convenience the JPM system does not require a valid username and
password combination for every privileged operation. To avoid this, JPM will
store the authentication token returned by the \registry. These are stored in a
file on the file-system.

TODO I can't write smart things right now. Words no good :-(

The authentication token (or session token) is saved instead of the username
and password.  Since if an attacker were to gain access to both the username
and password, she would be able to gain full control over the user, this
includes being able to generate new authentication tokens.

If an attacker instead were to gain only the authentication token, the damage
is not as bad. The attacker would still be able to pose as the user, but only
while the authentication token is valid.

Several steps are taken to mitigate the damage that can be done in case of
session hi-jacking. TODO I just called a session, without defining it.

% TODO Not sure if this needs to be implemented, but be careful when
% discussing. Don't make it sound like something is implemented if it isn't.
% Should make clear if something is used or not.

% We partially do this. It is possible, but needs to be implemented on the
% registry side of things. Would be an easy change.

Session timeout

% We do this

Manual session invalidation (logging out)

% We support this. Haven't defined any actions to be critical though.

Re-authentication before critical actions


\url{https://www.owasp.org/index.php/Authentication_Cheat_Sheet}

\url{https://www.owasp.org/images/0/08/OWASP_SCP_Quick_Reference_Guide_v2.pdf}

\section{Putting it all Together: Installing Dependencies}

In this section we will illustrate how the different services of JPM
communicates when the \txtl{install} command is invoked on a package. Most of
the individual components of this process has already been discussed, when
these are encountered we'll make a reference to the relevant sections.

TODO Finish this

\section{Examples and Discussion}

\subsection{Creating, Sharing, Using, and Upgrading a JPM Package}

In this example we will discover the workflow of creating a package, all the
way to using this package in another project.

First we will create a new JPM package by using the \txtl{init} command.
This will display prompts which guides the user through creating a new package.
Following this a new directory is created containing the basic files needed for
a package.

\begin{minted}{text}
dan@host:/ # jpm init
Package name
------------
> my-package

( ... Remaining cut for brevity ... )

dan@host:/ # cd my-package

dan@host:/my-package # ls
package.json

dan@host:/my-package # cat package.json | json
{
  "private": false,
  "name": "my-package",
  "description": "description",
  "version": "0.1.0",
  "authors": "Dan"
}
\end{minted}

We're now ready to actually implement our service. Here we'll create a simple
service exposing an interface called \joliel{MyPackageInterface} exposing a
single operation \joliel{hello(void)(string)} which will respond back with a
greeting.

With this service created, we can now publish the package to the public JPM
registry. All JPM commands which interact with a registry optionally takes a
registry name, if this is omitted it will default to the public registry.
Before we can publish a package we must create a user on the public registry.

\begin{minted}{text}
dan@host:/my-package # jpm register
Username
--------
> dan
Password:
Password (Repeat):
\end{minted}

It's possible to ask a registry who we are currently authenticated as:

\begin{minted}{text}
dan@host:/my-package # jpm whoami
dan
\end{minted}

When authenticated, it is possible to publish package to the public JPM
registry.

\begin{minted}{text}
dan@host:/my-package # jpm publish
\end{minted}

Now lets pretend we another user, who wishes to use our package. In order to
download from the JPM public registry there is no need to be authenticated
(although other registers may require this). Lets create a package and add it
as a dependency. This requires us to update the \txtl{package.json} files
generated by JPM:

\begin{minted}{json}
{
    "dependencies": [{ "name": "my-package", "version": "0.1.X" }]
}
\end{minted}

The dependency here used the SemVer expressions explained in Section
\ref{sec:versions}. This expression states that the newest version of the
package should be picked, as long as that version has a major version number of
0 and a minor version number of 1.

We can now download all the dependencies of our package:

\begin{minted}{text}
dan@host:/client # jpm install
Downloading          my-package@0.1.0
Completed            my-package@0.1.0

dan@host:/client # tree .
.
|-- jpm_lock.json
|-- jpm_packages
|   `-- my-package
|       |-- interface.iol
|       |-- main.ol
|       `-- package.json
`-- package.json

2 directories, 5 files
\end{minted}

We can see that this generates a lockfile which states that ``my-package'' has
been resolved from the public registry as version 0.1.0.

\begin{minted}{text}
dan@host:/client # cat jpm_lock.json | json
{
  "locked": {
    "my-package@0.1.X/public": {
      "resolved": "0.1.0"
    }
  },
  "_note": "Auto-generated"
}
\end{minted}

If we publish a new version of ``my-package'' and run \txtl{install}, then
we still be left with the same version. It is not until we run \txtl{upgrade}
that we receive the newest version.

\begin{minted}{text}
dan@host:/client # jpm search my-package
my-package@0.1.1/public
  description

dan@host:/client # jpm install
Downloading          my-package@0.1.0
Completed            my-package@0.1.0

dan@host:/client # jpm upgrade

dan@host:/client # jpm install
Downloading          my-package@0.1.1
Completed            my-package@0.1.1
\end{minted}

We can start our package by setting the ``main'' attribute of the package
manifest. This allows us to write:

\begin{minted}{text}
dan@host:/client # jpm run
Service running...
\end{minted}

We can even run with a particular configuration:

\begin{minted}{text}
dan@host:/client # jpm run --conf my-profile config.col
Service running...
\end{minted}

\subsection{The Calculator System}

Continuing with the ``modularized'' version of the calculator system from
Section \ref{sec:module_examples}.

The first step in making these usable from the package manager is of course to
create package manifests. This action becomes fairly straight forward,
especially since we already have a system working with the Jolie module system.
The package manifest for the \txtl{multiplication} package is shown in Listing
\ref{lst:multiplication_package}. The only significant multiplications here are
perhaps a version number, and a formalization of dependencies (in this case the
\txtl{numbers@1.0.0} package).

\begin{listing}[H]
\begin{minted}{json}
{
    "name": "multiplication",
    "description": "An multiplication package",
    "license": "MIT",
    "authors": ["Dan Sebastian Thrane <dathr12@student.sdu.dk>"],
    "main": "main.ol",
    "version": "1.0.0",
    "dependencies": [{ "name": "numbers", "version": "1.0.0" }]
}
\end{minted}
\caption{Package manifest for the \txtl{multiplication} package}
\label{lst:multiplication_package}
\end{listing}

The other package manifests will look fairly similar. The dependency graph ends
up mimicking the system architecture exactly. However there are some subtle
problems associated with the current approach.

Currently the calculator will always download the \txtl{numbers} package,
despite not directly depending on it. This happens due to the dependency on
\txtl{multiplication} which itself depends on \txtl{numbers}. However the
\txtl{calculator} service really only depends on \txtl{numbers} if
\txtl{multiplication} is embedded. At this point this may seem like a minor
detail, but this can get vastly more complicated if there are more dependencies
(which themselves may have dependencies and so on).

The solution to this problem, is to make sure the package manager will only
download the files we actually need, and no more. This would mean if no
embedding is ever performed, we will need only the dependencies required to
perform interfacing. If embedding is required we will need both the interface
dependencies along with a concrete implementation.

For the sake of simplicity the same dependency system is used. It is
recommended that every service separately publishes its interface and concrete
implementation. This approach has the added benefit of allowing multiple
implementations of same service. For closed source services a single package
containing only the interface could be published.

Depending on if an embedding is desired the concrete implementation can simply
be listed as a dependency, the interface dependency will implicitly be picked
up (from the concrete implementation's dependencies). If an embedding is
\emph{never} desired a dependency can be made to just the interface package.

This then begs the question of when to use a concrete implementation over a
dependency. TODO say something.

Additionally JPM might benefit from optional dependencies which are only when
certain conditions are made, for example during some ``development build''.
This might prove to be a decent compromise, allowing for the convenience of
embedding without directly forcing users of the package to also depend on it.

\subsection{Using Lifetime Hooks to Improve Development Workflow}

In this example we will discover how the use of Jolie lifetime hooks can
improve the development workflow. In this example we will look at a hybrid
Jolie-Java service. The service itself will be written in Jolie, using an
internal Java service to perform some of the computational heavy-lifting.  The
\txtl{packages} service, which performs package manifest validation, used in
the package manager is an example of such a package type. The primary flow and
communication is performed entirely by Jolie, while the more ``computationally
heavy'' parts, like parsing, are handled by the internal Java service.  This
architecture is illustrated in Figure \ref{fig:jolie_java_service}.

\begin{listing}[H]
\begin{minted}{text}
+----------------------------------------+
|                                        |
| jolie-service ---------+               |
|                        | local         |
|                        | communication |
|                        |               |
|                        v               |
|                      +----------------+|
|                      | java-service   ||
|                      +----------------+|
+----------------------------------------+
\end{minted}
\caption{A Jolie service using an internal Java service}
\label{fig:jolie_java_service}
\end{listing}

The Java service uses Gradle as its build system. The build system is capable
of producing a JAR file. The Jolie engine requires this file to be included in
the classpath in order to perform the embedding. If the file is placed in the
\txtl{lib} directory this will happen automatically.

A small bash script was made to run JAR creation through Gradle, and then copy
the artifact into the \txtl{lib} directory. Interested readers may find both
the Gradle build script and bash script in Appendix 3. Both scripts are
relatively simple, but should be a sufficient starting point.

The file structure of the project is shown in Listing
\ref{lst:java_jolie_file_structure}. The \txtl{build} file is the script which
performs the entire build and copies it into the \txtl{lib} directory.

\begin{listing}[H]
\begin{minted}{text}
.
|-- build
|-- java-service
|   |-- build.gradle
|   `-- src/main/java/dk/thrane/jolie/JavaService.java
|-- main.ol
`-- package.json
\end{minted}
\caption{}
\label{lst:java_jolie_file_structure}
\end{listing}

Assuming a lot of changes are made to the Java service, this would mean two
commands would almost always be required to create a new iteration of the
build. This can be minimized by creating a \txtl{pre-start} hook as shown in
Listing \ref{lst:pre_start_hook}. Having it as a part package manifest makes
the intent very clear. Having some home-made build script would instead rely
entirely on convention. Such conventions may vary from project to project,
while having it as a part of the package manifest would remain consistent
across all JPM packages.

\begin{listing}[H]
\begin{minted}{json}
{ "events": { "pre-start": "./build" } }
\end{minted}
\caption{A \txtl{pre-start} hook for performing compilation of the internal
    Java-service}
\label{lst:pre_start_hook}
\end{listing}

The JPM lifetime hooks aim to enforce the implicit workflow a project typically
has by making it explicit. This could for example be used to enforce the
constraint that testing must succeed before publishing to a registry. This can
be implemented by using the \txtl{pre-publish} hook. If the script being run
returns a non-zero exist code, then the publishing action will be canceled
(similar behavior applies to all \txtl{pre-*} hooks).

Currently however it would be useful if a Jolie package could allow for
multiple entry-points. This would make it easier to write tests written
directly in Jolie. However having multiple entry-points would bring up the
question of which to use for embedding. We could keep the constraint of a
single main entry-point used for embedding, and multiple others used for
various other tasks, such as testing.



\appendix
\chapter{Appendix}
\section{Appendix A: JPM Manifest Specification}\label{package-specification}

This document covers the specification of the file which defines a
package. The format used for this document will be JSON, but the format
and whether or not to allow for several documents is still up for
discussion. For now we should avoid using any features which the generic
Jolie value cannot support.

\subsection{Purpose}\label{purpose}

The purpose of the package document is to define what a package is.
Every Jolie package will contain such a document, and it describes
several important properties about the package. These properties are
described in the section ``Format and Properties''.

\subsection{Table of Contents}\label{table-of-contents}

\begin{itemize}
\tightlist
\item
  \protect\hyperlink{format-and-properties}{Format and Properties}

  \begin{itemize}
  \tightlist
  \item
    \protect\hyperlink{name}{name}
  \item
    \protect\hyperlink{version}{version}
  \item
    \protect\hyperlink{license}{license}
  \item
    \protect\hyperlink{authors}{authors}
  \item
    \protect\hyperlink{private}{private}
  \item
    \protect\hyperlink{main}{main}
  \item
    \protect\hyperlink{dependencies}{dependencies}
  \item
    \protect\hyperlink{dependency}{dependency}

    \begin{itemize}
    \tightlist
    \item
      \protect\hyperlink{name-1}{name}
    \item
      \protect\hyperlink{version-1}{version}
    \item
      \protect\hyperlink{registry}{registry}
    \end{itemize}
  \item
    \protect\hyperlink{registries}{registries}
  \item
    \protect\hyperlink{registry-1}{registry}

    \begin{itemize}
    \tightlist
    \item
      \protect\hyperlink{name-2}{name}
    \item
      \protect\hyperlink{location}{location}
    \end{itemize}
  \end{itemize}
\end{itemize}

\hypertarget{format-and-properties}{\subsection{Format and
Properties}\label{format-and-properties}}

\hypertarget{name}{\subsubsection{name}\label{name}}

\textbf{Name:} \texttt{name}

\textbf{Optional:} false

\textbf{Type:} \texttt{string}

\textbf{Description:} The \texttt{name} property uniquely defines a
package in a registry. Every registry must only contain a single package
with a given name.

\textbf{Rules:}

\begin{itemize}
\tightlist
\item
  The name of a package is \emph{not} case-sensitive
\item
  The length of a name is less than 255 characters
\item
  Names are US-ASCII
\item
  Names may only contain unreserved URI characters (see section 2.3 of
  \href{https://www.ietf.org/rfc/rfc3986.txt}{RFC 3986})
\end{itemize}

If any of these rules are broken the JPM tool should complain when
\emph{any} command is invoked. Similarly a registry should reject any
such package.

\hypertarget{version}{\subsubsection{version}\label{version}}

\textbf{Name:} \texttt{version}

\textbf{Optional:} false

\textbf{Type:} \texttt{string}

\textbf{Description:} This property describes the current version of
this package.

\textbf{Rules:}

\begin{itemize}
\tightlist
\item
  The version string must be a valid SemVer 2.0.0 string (see
  http://semver.org/spec/v2.0.0.html)
\end{itemize}

\hypertarget{license}{\subsubsection{license}\label{license}}

\textbf{Name:} \texttt{property\_name}

\textbf{Optional:} false

\textbf{Type:} \texttt{string}

\textbf{Description:} Describes the license that this package is under.

\textbf{Rules:}

\begin{itemize}
\tightlist
\item
  Must be a valid identifier. See https://spdx.org/licenses/
\end{itemize}

\hypertarget{authors}{\subsubsection{authors}\label{authors}}

\textbf{Name:} \texttt{authors}

\textbf{Optional:} false

\textbf{Type:}
\texttt{string\textbar{}array\textless{}string\textgreater{}}

\textbf{Description:} Describes the authors of this package

\textbf{Rules:}

\begin{itemize}
\tightlist
\item
  The array must contain at least a single entry
\item
  Each entry should follow this grammar:
\end{itemize}

\begin{verbatim}
name ["<" email ">"] ["(" homepage ")"]
\end{verbatim}

\hypertarget{private}{\subsubsection{private}\label{private}}

\textbf{Name:} \texttt{private}

\textbf{Optional:} true

\textbf{Type:} \texttt{boolean}

\textbf{Description:} Describes if this package should be considered
private. If a package is private it cannot be published to the
``public'' repository.

\textbf{Rules:}

\begin{itemize}
\tightlist
\item
  By default this property has the value of \texttt{true} to avoid
  accidential publishing of private packages.
\end{itemize}

\hypertarget{main}{\subsubsection{main}\label{main}}

\textbf{Name:} \texttt{main}

\textbf{Optional:} true

\textbf{Type:} \texttt{string}

\textbf{Description:} Describes the main file of a package.

\textbf{Rules:}

\begin{itemize}
\tightlist
\item
  The value is considered to be a relative file path from the package
  root.
\end{itemize}

\hypertarget{dependencies}{\subsubsection{dependencies}\label{dependencies}}

\textbf{Name:} \texttt{dependencies}

\textbf{Optional:} true

\textbf{Type:} \texttt{array\textless{}dependency\textgreater{}}

\textbf{Description:} Contains an array of dependencies. See the
``dependency'' sub-section for more details.

\textbf{Rules:}

\begin{itemize}
\tightlist
\item
  If the property is not listed, a default value of an empty array
  should be used
\end{itemize}

\hypertarget{dependency}{\subsubsection{dependency}\label{dependency}}

\textbf{Type:} \texttt{object}

\textbf{Description:} A dependency describes a single dependency of a
package. This points to a package at a specific point on a specific
registry.

\hypertarget{name-1}{\paragraph{name}\label{name-1}}

\textbf{Name:} \texttt{name}

\textbf{Optional:} false

\textbf{Type:} \texttt{string}

\textbf{Description:} Describes the name of the dependency. This refers
to the package name, as defined earlier.

\textbf{Rules:} A dependency name follows the exact same rules as a
package name.

\hypertarget{version-1}{\paragraph{version}\label{version-1}}

\textbf{Name:} \texttt{version}

\textbf{Optional:} false

\textbf{Type:} \texttt{string}

\textbf{Description:} Describes the version to use

\textbf{Rules:}

\begin{itemize}
\tightlist
\item
  Must be a valid SemVer 2.0.0 string
\item
  (This property follows the same rules as the package version does)
\end{itemize}

\hypertarget{registry}{\paragraph{registry}\label{registry}}

\textbf{Name:} \texttt{registry}

\textbf{Optional:} true

\textbf{Type:} \texttt{string}

\textbf{Description:} This describes the exact registry to use. If no
registry is listed the ``public'' registry will be used.

\textbf{Rules:}

\begin{itemize}
\tightlist
\item
  The value of this property must be a valid registry as listed in the
  \texttt{registries} property.
\end{itemize}

\hypertarget{registries}{\subsubsection{registries}\label{registries}}

\textbf{Name:} \texttt{registries}

\textbf{Optional:} true

\textbf{Type:} \texttt{array\textless{}registry\textgreater{}}

\textbf{Description:} Contains an array of known registries. See the
registry sub-section for more details.

\textbf{Rules:}

\begin{itemize}
\tightlist
\item
  This property contains an implicit entry which points to the public
  registry. This registry is named ``public''.
\end{itemize}

\hypertarget{registry-1}{\subsubsection{registry}\label{registry-1}}

\textbf{Type:} \texttt{object}

\textbf{Description:} A registry describes a single JPM registry. A JPM
registry is where the package manager can locate a package, and also
request a specific version of a package.

\hypertarget{name-2}{\paragraph{name}\label{name-2}}

\textbf{Name:} \texttt{name}

\textbf{Optional:} false

\textbf{Type:} \texttt{string}

\textbf{Description:} This property uniquely identifies the registry.

\textbf{Rules:}

\begin{itemize}
\tightlist
\item
  A name cannot be longer than 1024 characters
\item
  The name cannot be ``public''
\item
  No two registries may have the same name
\end{itemize}

\textbf{TODO:}

\begin{itemize}
\tightlist
\item
  Encoding of name
\item
  Should the length limit be dropped? There is no technical reason for
  the limit
\end{itemize}

\hypertarget{location}{\paragraph{location}\label{location}}

\textbf{Name:} \texttt{location}

\textbf{Optional:} false

\textbf{Type:} \texttt{string}

\textbf{Description:} Describes the location of the registry.

\textbf{Rules:}

\begin{itemize}
\tightlist
\item
  Must be a valid Jolie location string (e.g.
  ``socket://localhost:8080'')
\end{itemize}

\section{COL Grammar}\label{package-specification}

Presented as an ABNF-like syntax. All literals are case sensitive.

\begin{minted}{abnf}
configuration-tree = *include *region

include = "include" qstring

region = region-header "{" region-body "}"
region-header = [ "profile" qstring ] "configures" qstring [ "extends" qstring ]
region-body = *definition
definition = port | interface | parameter

port = input-port | output-port
input-port = "inputPort" identifier "{" port-body "}"
output-port = embedded-output-port | external-output-port
external-output-port = "outputPort" identifier "{" port-body "}"
embedded-output-port = "outputPort" identifier "embeds" qstring "with" qstring
port-body = *port-property
port-property = location-property | protocol-property
location-property = "Location" ":" qstring
protocol-property = "Protocol" ":" identifier [ protocol-config ]
protocol-config = inline-tree

interface = "interface" identifier "=" identifier "from" qstring

parameter = variable-path "=" value
variable-path = var-id *var-node
var-id = ( "(" qstring ")" ) | identifier
var-node = "." var-id [ "[" unsigned-int "]"  ]
value = ( primitive [ inline-tree ] ) | inline-tree
inline-tree = "{" *(tree-child ",") tree-child "}"
tree-child = "." variable-path "=" value
primitive = qstring | int | long | double | bool

letter = <any Unicode code point recognized by Java via Character#isLetter>
digit = "0" | "1" | "2" | "3" | "4" | "5" | "6" | "7" | "8" | "9"
bool = "true" | "false"
int = unsigned-int | "-" unsigned-int
long = int ( "l" | "L" )
unsigned-int = 1*digit ; Jolie discards leading 0s
double = 1*digit "." 1*digit [ double-exp ] | 1*digit double-exp
double-exp = ( "e" | "E" ) [ "-" | "+" ] 1*digit
qstring  = ( <"> *(qdtext | quoted-pair ) <"> )
quoted-pair = "\" ( "\" | "n" | "t" | "r" | <"> | "u" )
qdtext = <any TEXT except <">>
\end{minted}



\printbibliography
\end{document}
